% TOP_PROBLEM: {"id": 6, "title": "Hodge Conjecture", "category_id": 5, "category": {"id": 5, "name": "algebraic_geometry", "display_name": "Algebraic Geometry", "description": "Geometric objects defined by polynomial equations.", "slug": "algebraic-geometry", "order_index": 5, "created_at": "2026-07-31T15:26:25.671Z"}, "status": "open"}
% FIRST_POSTED: 2026-09-27
% ATTEMPT_STATUS: unresolved
% Consolidated research record by GPT-6 Astra Pro, 27 September 2026.
% Earlier mathematical arguments and finite certificates are preserved with their limitations.
% Compile with LuaLaTeX or XeLaTeX. No external bibliography or image files are required.
% The four original research ZIP archives are embedded as comments after \end{document}.
\documentclass[11pt]{article}
\usepackage[margin=25mm]{geometry}
\usepackage{fontspec}
\setmainfont{Latin Modern Roman}
\usepackage{amsmath,amssymb,mathtools}
\usepackage{unicode-math}
\setmathfont{Latin Modern Math}
\usepackage{longtable,booktabs,array}
\usepackage{fvextra}
\usepackage{xurl,hyperref}
\hypersetup{colorlinks=true,urlcolor=blue,linkcolor=blue,
  pdftitle={Hodge Conjecture: Complete Research Record},
  pdfauthor={GPT-6 Astra Pro},bookmarksdepth=2}
\setcounter{secnumdepth}{0}
\setlength{\parindent}{0pt}
\setlength{\parskip}{5pt}
\setlength{\emergencystretch}{3em}
\setlength{\LTpre}{4pt}
\setlength{\LTpost}{5pt}
\renewcommand{\arraystretch}{1.13}
\newcolumntype{L}[1]{>{\raggedright\arraybackslash}p{#1}}
\allowdisplaybreaks[1]
\newcommand{\Q}{\mathbb Q}
\newcommand{\Z}{\mathbb Z}
\newcommand{\C}{\mathbb C}
\newcommand{\F}{\mathbb F}
\newcommand{\cl}{\operatorname{cl}}
\newcommand{\CH}{\operatorname{CH}}
\newcommand{\Res}{\operatorname{Res}}
\newcommand{\PD}{\operatorname{PD}}
\newcommand{\Tr}{\operatorname{Tr}}
\newcommand{\src}[1]{\hyperlink{src:#1}{[#1]}}
\DefineVerbatimEnvironment{Code}{Verbatim}{fontsize=\small,breaklines=true,breakanywhere=true}
\begin{document}
% problemId: 6
% Canonical problem ID: 6
\section*{\texorpdfstring{Hodge Conjecture}{Hodge Conjecture}}\label{hodge}\label{6}
\subsection{Definitions and mathematical statement}
Let $X$ be a smooth projective algebraic variety over $\C$ of complex dimension
$n$. Its singular cohomology has the Hodge decomposition
\[
 H^k(X,\C)=\bigoplus_{r+s=k}H^{r,s}(X),\qquad
 \overline{H^{r,s}(X)}=H^{s,r}(X),
\]
and Hodge filtration
\[
 F^pH^k(X,\C)=\bigoplus_{r\ge p}H^{r,k-r}(X).
\]
Here $H^{r,s}(X)$ is the component represented by closed forms of bidegree
$(r,s)$, equivalently by harmonic forms of that type after a K\"ahler metric
has been chosen. A rational Hodge class of codimension $p$ is an element of
\[
 \operatorname{Hdg}^{2p}(X)
 :=H^{2p}(X,\Q)\cap H^{p,p}(X)
 =H^{2p}(X,\Q)\cap F^pH^{2p}(X,\C),
\]
where rational cohomology is identified with its image in complex cohomology.
The equality uses the reality of a rational class. These conventions are those
of Deligne's supplied problem description, Section 1.\src{E1}

If $Z\subset X$ is an algebraic subvariety of complex codimension $p$, its
fundamental cycle has Poincar\'e dual
$\cl(Z)\in H^{2p}(X,\Z)$, whose complex image lies in $H^{p,p}(X)$.
An algebraic cycle is a finite integral linear combination of such
subvarieties. Write $\CH^p(X)$ for codimension-$p$ cycles modulo rational
equivalence; the cycle-class map factors through this quotient. With
$\CH^p(X)_{\Q}:=\CH^p(X)\otimes_{\Z}\Q$, the exact selected assertion is
\begin{equation}\label{hodge:statement}
 \boxed{\quad
 \cl\bigl(\CH^p(X)_{\Q}\bigr)=\operatorname{Hdg}^{2p}(X)
 \quad\text{for every smooth projective }X/\C\text{ and }0\le p\le n.
 \quad}
\end{equation}
Equivalently, for every rational Hodge class $\alpha$ there are finitely many
codimension-$p$ algebraic subvarieties $Z_j$ and coefficients $q_j\in\Q$ such
that $\alpha=\sum_jq_j\cl(Z_j)$. The counterexample requested at the start of
the conversation would therefore have to satisfy
\begin{equation}\label{hodge:counterexample}
 X/\C\text{ smooth projective},\qquad
 \alpha\in\operatorname{Hdg}^{2p}(X),\qquad
 \alpha\notin\cl\bigl(\CH^p(X)_{\Q}\bigr).
\end{equation}
An integral obstruction that disappears after tensoring with $\Q$, or a
counterexample on a non-projective compact K\"ahler manifold, does not answer
this question. Deligne explicitly discusses both distinctions in Section 2
of the supplied text.\src{E1}

\par\smallskip\textit{Formulation and concept references:}
\href{https://www.claymath.org/library/monographs/MPPc.pdf}{[S1]},
\href{https://www.claymath.org/wp-content/uploads/2022/06/hodge.pdf}{[E1]}.
The uploaded \texttt{hodge.pdf}, rather than a reported solution or a stronger
variant, is the authoritative problem statement for this record.

\subsection{Short English statement}
On a smooth complex projective variety, must every rational cohomology class
with the Hodge type of an algebraic subvariety be a rational linear combination
of classes of actual algebraic subvarieties?

\subsection{Research attempt}
\paragraph{Scope, chronology, and evidential status.}
This document consolidates the complete mathematical progression in the
conversation and the four supplied research archives. It follows the example's
organization: definitions, a short English statement, a research attempt, and
sources. The example concerns BSD; its mathematical content and catalogue rank
are not transferred to Hodge.\src{T1}

The initial instruction was to search for a counterexample, not to prove the
conjecture. Subsequent requests asked for progressively stronger attempts,
including a possible response to rumors of an unpublished AI-assisted
breakthrough. The last mathematical response reversed the local search and
asked for a surface that would eliminate the chosen candidate as a
counterexample. That reversal is recorded as a separate strategy, not
silently substituted for the original task.

\textbf{Status of the record.} No verified counterexample, general proof,
or algebraic surface meeting the final witness conditions was obtained.
The finite computations below were reported in the earlier rounds, and their
source code and outputs are retained. The geometric arguments are inherited
research arguments, not independently peer-reviewed or Lean-formalized
results. Historical novelty was not established. Exporting them does not
upgrade their evidential status. Where an established theorem is invoked, it
is identified separately from the specialization made in the conversation.

There are four computational stages: the initial Fermat-character search;
the arithmetic and deformation continuation; the second-order two-monomial
obstruction; and the explicit-period and finite-cover pass. Two later
conversational stages concern the reported breakthrough and the
complete-intersection witness strategy. Their mathematical and provenance
content is retained below. The news discussion is a record of what the
conversation asserted, not a new verification of events as of 27 September
2026.\src{C1}\src{C2}\src{C3}\src{C4}\src{C5}

\paragraph{Round I: the concrete smooth projective candidate.}
The search focused on codimension two on the degree-$70$ Fermat fourfold
\begin{equation}\label{hodge:fermat}
 X=X_0=\{F_0=0\}\subset\mathbb P^5_{\C},\qquad
 F_0=x_0^{70}+x_1^{70}+x_2^{70}+x_3^{70}+x_4^{70}+x_5^{70},
\end{equation}
and the ordered character
\begin{equation}\label{hodge:character}
 a=(1,20,24,42,61,62).
\end{equation}
Projectivity is immediate. The derivatives are $70x_i^{69}$, so their
simultaneous vanishing would force all homogeneous coordinates to be zero;
therefore $X$ is smooth.

Let $U=(\Z/70\Z)^\times$, let $\zeta=\exp(2\pi i/70)$, and set
\[
 b(t)=(\langle ta_0\rangle_{70},\ldots,
       \langle ta_5\rangle_{70})\quad(t\in U),
 \qquad B=\{b(t):t\in U\},
\]
where $\langle\cdot\rangle_{70}$ is the least positive residue. The diagonal
coordinate-scaling group is $(\mu_{70})^6/\mu_{70}$, with the denominator
acting by common projective scaling. Since $\sum_i a_i=210\equiv0\pmod{70}$,
$a$ defines a character of this quotient. The standard Fermat decomposition
assigns a one-dimensional primitive cohomology eigenspace $V_b$ to each
relevant ordered character. The rational Hodge criterion for a fourfold is
\begin{equation}\label{hodge:balance}
 \sum_{i=0}^{5}\langle ta_i\rangle_{70}=3\cdot70=210
 \qquad\text{for every }t\in U.
\end{equation}
Requiring every Galois conjugate, rather than one complex eigenline, is
essential to the rational statement.\src{E2}

The recorded checks verify \eqref{hodge:balance} for all $24$ units. Since
$a_0=1$, the $24$ ordered characters are distinct. Thus the rational
character block $W$ specified by
\begin{equation}\label{hodge:block}
 W\otimes_{\Q}\C=\bigoplus_{t\in U}V_{b(t)}
\end{equation}
has dimension $24$ and lies in $H^{2,2}(X)$. At this first stage, $W$ was a
specified subspace, not a distinguished rational vector and not a
non-algebraicity certificate.\src{C1}

For completeness, the complete ordered residue check from the saved output is
reproduced here. It should not be confused with the coarser enumeration modulo
coordinate permutations used later.

\begin{longtable}{rcr}
\toprule
$t$ & $b(t)$ & Sum\\
\midrule
\endhead
$1$ & $(1,20,24,42,61,62)$ & $210$\\
$3$ & $(3,60,2,56,43,46)$ & $210$\\
$9$ & $(9,40,6,28,59,68)$ & $210$\\
$11$ & $(11,10,54,42,41,52)$ & $210$\\
$13$ & $(13,50,32,56,23,36)$ & $210$\\
$17$ & $(17,60,58,14,57,4)$ & $210$\\
$19$ & $(19,30,36,28,39,58)$ & $210$\\
$23$ & $(23,40,62,56,3,26)$ & $210$\\
$27$ & $(27,50,18,14,37,64)$ & $210$\\
$29$ & $(29,20,66,28,19,48)$ & $210$\\
$31$ & $(31,60,44,42,1,32)$ & $210$\\
$33$ & $(33,30,22,56,53,16)$ & $210$\\
$37$ & $(37,40,48,14,17,54)$ & $210$\\
$39$ & $(39,10,26,28,69,38)$ & $210$\\
$41$ & $(41,50,4,42,51,22)$ & $210$\\
$43$ & $(43,20,52,56,33,6)$ & $210$\\
$47$ & $(47,30,8,14,67,44)$ & $210$\\
$51$ & $(51,40,34,42,31,12)$ & $210$\\
$53$ & $(53,10,12,56,13,66)$ & $210$\\
$57$ & $(57,20,38,14,47,34)$ & $210$\\
$59$ & $(59,60,16,28,29,18)$ & $210$\\
$61$ & $(61,30,64,42,11,2)$ & $210$\\
$67$ & $(67,10,68,14,27,24)$ & $210$\\
$69$ & $(69,50,46,28,9,8)$ & $210$\\
\bottomrule
\end{longtable}

\paragraph{Round I: exhaustive sextuple search and its six residual cases.}
The first program enumerated all nondecreasing sextuples with entries in
$\{1,\ldots,69\}$ and total $210$, tested the Hodge condition, and then
identified cases under coordinate permutation and multiplication by a unit.
An independent Python method matched complete residue-sum signatures of
triples; dynamic programming separately checked the total number of
sum-$210$ sextuples. The saved results are as follows.\src{C1}
\begin{center}
\begin{tabular}{lr}
\toprule
Enumeration or successive screen & Number\\
\midrule
Nondecreasing sextuples of sum $210$ & $1,426,456$\\
Sextuples satisfying all rational Hodge tests & $10,948$\\
Cases modulo permutations and units & $882$\\
Containing a complementary pair & $843$\\
Containing a zero-sum triple after the first screen & $6$\\
Decomposable after adjoining one complementary pair & $27$\\
Residual after these three successive screens & $6$\\
\bottomrule
\end{tabular}
\end{center}
The last four rows partition the $882$ cases in the indicated order. The
third construction screen adjoins a pair $(r,70-r)$ and tests whether the
resulting eight entries split into two Hodge quadruples. The six residual
representatives saved in \texttt{residual70.txt} are
\begin{align*}
 &(1,6,43,45,57,58), &&(1,15,29,43,57,65),\\
 &(1,20,24,42,61,62),&&(1,20,29,43,57,60),\\
 &(1,29,35,43,45,57),&&(2,16,30,44,58,60).
\end{align*}
The chosen character has no pair summing to $70$, no triple summing to zero
modulo $70$, and no successful one-pair decomposition into Hodge quadruples.
The earlier conversation identified degree $70$ as a search direction from a
reported residual-case computation in the preprint cited as \src{E4}; that
attribution is inherited, not a new validation of that preprint.

\paragraph{Round I: the formal character lattice is not the cycle-class lattice.}
The scripts assign to a tuple $c$ its multiplicity vector
\[
 \nu(c)=\sum_{j}\mathbf e_{c_j}\in\Z^{m-1},
\]
where $\mathbf e_r$ records entry $r\in\{1,\ldots,m-1\}$. The tested lattice
$S_m$ is generated by the vectors of the complementary pairs $(r,m-r)$ and
the following specified standard tuples, with $p$ a prime divisor of $m$ and
$1\le r<m/p$:
\[
 \begin{cases}
 (r,r+m/2,m-2r,m/2),&p=2,\\
 (r,r+m/p,\ldots,r+(p-1)m/p,m-pr),&p\text{ odd}.
 \end{cases}
\]
These are the generator definitions actually used in the supplied programs;
the geometric context is attributed there to Aoki and the Fermat
construction literature.\src{E2}\src{E3}\src{E4}\src{C1}

Exact Hermite-normal-form calculations at $m=70$ give lattice rank $57$ and
\begin{equation}\label{hodge:formal-lattice}
 \nu(a)\notin S_{70},\qquad 2\nu(a)\in S_{70}.
\end{equation}
Of the six residual representatives, the selected $a$ is the only one outside
this specified lattice. The coordinate denominators in a lattice basis are
$1$ and $2$. This is not an integral-versus-rational obstruction for a
cohomology class on $X$: formal concatenation or multiplicity of character
tuples is not scalar multiplication of a cycle class on the same fourfold.
In particular, \eqref{hodge:formal-lattice} says nothing by itself about
whether $\alpha$ or $2\alpha$ is algebraic for $\alpha\in W$.

\paragraph{Round I: Jacobi sums and the failed non-root-of-unity test.}
For a split prime $p\equiv1\pmod{70}$, choose a primitive root $g_p$ and the
multiplicative character $\chi_p$ with $\chi_p(g_p)=\zeta$. The six-character
Jacobi sum at additive total one is computed in the exact cyclotomic ring
$\Z[\zeta]\simeq\Z[T]/(\Phi_{70}(T))$. In the saved convention, the
fourfold Frobenius eigenvalue is $\lambda_p=-J_p$, equivalently the
zero-additive-total sum divided by $p-1$. The relation of diagonal
hypersurface point counts, Jacobi sums, and Frobenius is an established input,
not a result of the finite search.\src{E5}\src{E6}

An algebraic codimension-two cycle defined over $\F_{p^r}$ has a class fixed
by the $r$th power of normalized Frobenius. Thus
\[
 \operatorname{Frob}_p^{\,r}\cl(Z)=p^{2r}\cl(Z),
\]
and a normalized eigenvalue $\lambda_p/p^2$ that is not a root of unity would
exclude an algebraic representative over every finite extension of that
finite field. This was the attempted arithmetic obstruction.

The first round tested $71$, $211$, and $281$; the continuation added $491$.
Full additive convolution and a product of binary Jacobi sums were reported
to agree exactly. The full four-prime record is
\begin{center}
\begin{tabular}{rrrr}
\toprule
$p$ & $g_p$ & $\lambda_p/p^2$ & Order\\
\midrule
$71$ & $7$ & $\zeta^{40}$ & $7$\\
$211$ & $2$ & $\zeta^{12}$ & $35$\\
$281$ & $3$ & $\zeta^{32}$ & $35$\\
$491$ & $2$ & $\zeta^{57}$ & $70$\\
\bottomrule
\end{tabular}
\end{center}
For example, at $491$ the additive-total-one sum is
$491^2\zeta^{22}$, whose negative is $491^2\zeta^{57}$. Changing the embedding
or Frobenius convention may conjugate or invert the displayed roots, without
changing their orders.\src{C1}\src{C2}

At $71$, all normalized eigenvalues on $W$ have order $7$, so a
codimension-two cycle defined over $K=\Q(\zeta)$ cannot have a nonzero
projection into this block. But cycles over finite extensions of $K$, and in
particular cycles over $\C$, are not thereby excluded. The first three
values all have $35$th power one; the fourth has exact order $70$. The
hoped-for non-root-of-unity obstruction did not occur.

\paragraph{Round II: why further primes cannot repair that arithmetic attack.}
The continuation invokes the Jacobi-sum Hecke-character description. The
constant residue sums make the normalized character have infinity type zero,
hence finite order. Its values lie in $K$, whose roots of unity are
$\mu_{70}$, so its order divides $70$. The value of order $70$ at $491$
therefore makes the global normalized character have exact order $70$.
This is a structural deduction from the Hecke-character description, not an
extrapolation from four numerical samples.\src{E6}\src{C2}

The recorded class-field-theoretic consequence is a cyclic extension
\begin{equation}\label{hodge:abstract-field}
 L/K,\qquad [L:K]=70,
\end{equation}
over which normalized Galois action on $W$ is trivial. The unit-power
conjugate characters have the same kernel. A field of definition supporting
a nonzero cycle component in $W$ must satisfy this kernel condition; in the
recorded split-prime situation, residue degrees above $491$ must kill an
order-$70$ Frobenius value. These are necessary field-of-definition
conditions, not sufficient conditions for algebraicity. The explicit
period-field calculation in Round IV gives a concrete radical matching this
arithmetic behavior.

\paragraph{Round II: exhaustive straight monomial deformations.}
The next strategy was to move off the Fermat equation while keeping a
rational Hodge class. Consider
\[
 F_s=F_0+s x^\beta,\qquad |\beta|=70,\qquad0\le\beta_i\le68.
\]
The inequalities choose the standard monomial basis modulo infinitesimal
coordinate changes. There are
\begin{equation}\label{hodge:monomial-count}
 \binom{75}{5}-36=17,259,354
\end{equation}
such directions. The omitted monomials are the six $x_i^{70}$ and the thirty
$x_i^{69}x_j$ with $i\ne j$, spanning the degree-$70$ Jacobian ideal.

The exact escape criterion seeks the least $j\ge1$ for which some $b\in B$
and $q\in\{1,2\}$ satisfy
\begin{equation}\label{hodge:escape}
 r_i\equiv b_i-j\beta_i\pmod{70},\qquad
 1\le r_i\le69,\qquad\sum_i r_i=70q.
\end{equation}
The test form $\Res(x^{r-\mathbf1}\Omega/F_s^q)$ belongs to $F^3$ and must
pair to zero with the flat transport of a $(2,2)$-class. Taylor expansion and
Griffiths--Dwork reduction give a coefficient proportional to the
$b$-character period. If
\[
 \ell_i=\frac{r_i+j\beta_i-b_i}{70},
\]
then all $\ell_i$ are nonnegative, $\sum_i\ell_i=q+j-3$, and the exact
factor used in the archived calculation is
\begin{equation}\label{hodge:reduction-factor}
 \frac{(-1)^j2}{j!}\prod_{i=0}^{5}
       \left(\frac{b_i}{70}\right)_{\ell_i}.
\end{equation}
Here $(z)_k=z(z+1)\cdots(z+k-1)$, with $(z)_0=1$. The factor is nonzero.
The residue pattern is periodic in $j$ modulo the order of $\beta$ in
$(\Z/70\Z)^6$, so the implemented criterion is finite. The pole reduction
identity is an established method; the enumeration and specialization are
from the supplied continuation.\src{E7}\src{C2}

For a nonzero rational vector in $W$, every ordered character component is
nonzero: otherwise Galois conjugacy of the projectors forces all components
to vanish. The particular Taylor coefficient in this single-monomial test
selects one character and has no competing component to cancel it. This is
why the recorded escape test applies to every nonzero rational class in
$W$, rather than just a selected complex eigenvector.

The saved census is
\begin{center}
\begin{tabular}{rr}
\toprule
First escape order & Number of directions\\
\midrule
$1$ & $17,244,705$\\
$2$ & $13,617$\\
$3$ & $862$\\
$4$ & $105$\\
$5$ & $35$\\
$6$ & $15$\\
$7$ & $0$\\
$8$ & $15$\\
No escape & $0$\\
\bottomrule
\end{tabular}
\end{center}
Thus $14,649$ directions pass the first-order screen, but every straight
monomial family fails by order eight. This statement does not exclude
curved deformations with nonlinear corrections, or cancellations between
distinct monomials for an individual class.\src{C2}

\paragraph{Round II: an eighth-order obstruction missed by a seventh-order truncation.}
For
\[
 \beta=(2,0,0,68,0,0),\qquad F_s=F_0+s x_0^2x_3^{68},
\]
the recorded Hodge-period tests vanish through order seven. At order eight
choose
\[
 b=b(23)=(23,40,62,56,3,26),\qquad
 r=(7,40,62,2,3,26),\qquad q=2.
\]
The initial numerator is
\[
 x^{r-\mathbf1}=x_0^6x_1^{39}x_2^{61}x_3x_4^2x_5^{25}.
\]
Multiplication by $(x^\beta)^8$ gives exponent $545$ on $x_3$;
seven reductions lower that exponent to $55$ and the pole from $10$ to $3$.
With $\ell=(0,0,0,7,0,0)$, the coefficient is
\begin{equation}\label{hodge:eighth}
 \frac{2}{8!}\left(\frac45\right)_7
 =\frac{56202}{390625}\ne0.
\end{equation}
Writing $I_b(0)$ for the nonzero target-character pairing gives
\[
 I_r(s)=\frac{56202}{390625}I_b(0)s^8+O(s^9).
\]
This is a failure of the proposed class to remain Hodge, not a counterexample
to the Hodge conjecture.\src{C2}

\paragraph{Round II: the simultaneous tangent space and the degree-$71$ certificate.}
Let
\[
 S=\C[x_0,\ldots,x_5],\quad
 J=(x_0^{69},\ldots,x_5^{69}),\quad R=S/J,
 \qquad P_b=x^{b-\mathbf1}\in R_{204}.
\]
The period-matrix tangent description gives the simultaneous tangent space
for all of $W$ as
\begin{equation}\label{hodge:tangent}
 T_W=\{h\in R_{70}:hP_b=0\text{ in }R_{274}\text{ for every }b\in B\}.
\end{equation}
The general tangent formula is an established input.\src{E8}
Since $B$ is closed under $b\mapsto70-b$, a monomial $x^\beta$ lies in this
tangent space exactly when, for every $b\in B$, at least one coordinate has
$\beta_i\ge b_i$.

The coordinatewise gcds are
\[
 (g_0,\ldots,g_5)
 =\bigl(\gcd(70,a_0),\ldots,\gcd(70,a_5)\bigr)
 =(1,10,2,14,1,2).
\]
Put
\[
 \mathfrak d=(x_0^{69},x_1^{60},x_2^{68},x_3^{56},x_4^{69},x_5^{68}).
\]
The archived degree-$70$ result is
\begin{equation}\label{hodge:tangent-ideal}
 T_W=(\mathfrak d/J)_{70}.
\end{equation}
It is not an assertion about equality of full ideals in every degree. An
independent covering computation finds the first additional annihilating
monomial outside $\mathfrak d$ at degree $71$, for example
\[
 x_1^{50}x_2^2x_4x_5^{18},\qquad
 \beta=(0,50,2,0,1,18).
\]
The meet-in-the-middle certificate uses $365$ compressed masks for the first
three coordinates and $279$ for the last three; joining them yields minimum
total degree $71$. Hence no extra annihilator occurs in degree $70$.
The six generator supports of $\mathfrak d_{70}$ are disjoint, because any
pair of their degrees sums to more than $70$. Therefore
\begin{equation}\label{hodge:tangent-dim}
 \dim\mathfrak d_{70}=\sum_{i=0}^{5}\binom{g_i+5}{5}=14,685,
 \qquad \dim T_W=14,685-36=14,649.
\end{equation}
A separate slab enumeration of precisely these $14,649$ monomials reproduces
the saved higher-order escape counts without traversing the entire original
search space.\src{C2}

\paragraph{Round II: nonlinear finite covers preserving the whole block.}
The continuation then considered homogeneous polynomials $Q_i$ of degrees
$g_i$, close to
\[
 (Q_0,Q_1,Q_2,Q_3,Q_4,Q_5)
 =(x_0,x_1^{10},x_2^2,x_3^{14},x_4,x_5^2),
\]
and the family
\begin{equation}\label{hodge:power-family}
 X_Q=\{Q_0^{70}+Q_1^7+Q_2^{35}+Q_3^5+Q_4^{70}+Q_5^{35}=0\}.
\end{equation}
Nearby members are smooth and map to the fixed weighted hypersurface
\begin{equation}\label{hodge:weighted}
 \begin{gathered}
 Y=\{y_0^{70}+y_1^7+y_2^{35}+y_3^5+y_4^{70}+y_5^{35}=0\},\\
 Y\subset\mathbb P(1,10,2,14,1,2),\qquad
 q_Q:[x]\longmapsto[Q_0(x):\cdots:Q_5(x)].
 \end{gathered}
\end{equation}
At the Fermat point this is the quotient by
$H=\mu_{10}\times\mu_2\times\mu_{14}\times\mu_2$, acting on coordinates
$1,2,3,5$. Its degree is $560$. Every character in $B$ is $H$-invariant,
so the recorded argument identifies $W$ as the pullback of a rational Hodge
subspace $W_Y\subset H^4(Y,\Q)$. The maps remain everywhere defined and
finite in a sufficiently small neighborhood, and pullback of the fixed
subspace gives the all-orders Hodge-preserving family.\src{C2}

For example,
\begin{equation}\label{hodge:curved-example}
 \sum_{i\ne3}x_i^{70}
 +\left(x_3^2+\frac{s}{35}x_0^2\right)^{35}=0
\end{equation}
has the same initial direction $x_0^2x_3^{68}$ as the straight family, but
its nonlinear corrections keep it in the power family. Thus the eighth-order
failure of the straight line is compatible with all-orders preservation
along a curved family.

\paragraph{Round II: the inherited local-completeness argument.}
Work in the open affine space of smooth degree-$70$ equations with a local
topological marking. The map
\[
 \Phi(Q)=\sum_i Q_i^{70/g_i}
\]
has derivative at the Fermat tuple
\[
 d\Phi(P)=\sum_i\frac{70}{g_i}x_i^{70-g_i}P_i,
 \qquad \deg P_i=g_i.
\]
The disjoint monomial supports make this derivative injective, with image
$\mathfrak d_{70}$ of dimension $14,685$. Its image is therefore a smooth
analytic germ contained in the simultaneous Hodge locus $V_W$. By
\eqref{hodge:tangent-ideal}, that locus has tangent dimension exactly
$14,685$ in affine equation space. The contained smooth germ gives the
matching lower bound on local dimension, so local dimension equals embedding
dimension. The recorded argument concludes that $V_W$ is regular, smooth,
and reduced at the Fermat point, and its germ equals the power-family germ.
A transverse marked slice has dimension $14,649$; this should not be confused
with a coarse moduli quotient by a finite automorphism group.\src{C2}

This argument is preserved as an inherited geometric deduction. The finite
covering and monomial computations are certificates for its stated tangent
dimensions, not a formal proof-assistant check of the analytic geometry. The
claim concerns preservation of the entire block, not the larger possible
Hodge locus of one rational vector. Later period calculations do not require
this local-completeness claim.

\paragraph{Round II: why the power family transports the same algebraicity question.}
The intermediate $Y$ may be singular; it is not itself proposed as a
counterexample to the smooth-projective statement. Choose a smooth
projective resolution $f:\widetilde Y\to Y$. For $\beta\in W_Y$, put
$\widetilde\beta=f^*\beta$ and $\alpha_Q=q_Q^*\beta$. Resolving the graph
of the rational map $X_Q\dashrightarrow\widetilde Y$ gives maps
\[
 u:\Gamma\longrightarrow X_Q,\qquad
 v:\Gamma\longrightarrow\widetilde Y,
\]
with $u$ birational, $v$ generically finite of degree $560$, and
\[
 u^*\alpha_Q=v^*\widetilde\beta,\qquad
 u_*u^*\alpha_Q=\alpha_Q,\qquad
 v_*v^*\widetilde\beta=560\widetilde\beta.
\]
Pullback and pushforward between the smooth projective varieties preserve
rational algebraic cycle classes. Thus the recorded correspondence argument
is
\begin{equation}\label{hodge:algebraicity-invariance}
 \alpha_Q\text{ is algebraic}
 \quad\Longleftrightarrow\quad
 \widetilde\beta\text{ is algebraic}.
\end{equation}
This neither proves nor disproves algebraicity. It explains why varying $Q$
inside this family cannot change its answer.\src{C2}

\paragraph{Round II: the proposed escape by preserving only one rational class.}
The next proposed direction allowed cancellation between two monomials.
Write a Jacobian-ring representative in the form
\[
 P_\alpha=\sum_{t\in U}c_t x^{b(t)-\mathbf1},
\]
with the coefficient convention refined by period coordinates below. The
coefficients of a rational class cannot be chosen as arbitrary complex
numbers. Take
\begin{equation}\label{hodge:beta-pair}
 \begin{aligned}
 \beta_1&=(3,11,1,15,23,17),\\
 \beta_2&=(13,1,31,15,3,7),\\
 h&=\lambda x^{\beta_1}+\mu x^{\beta_2}.
 \end{aligned}
\end{equation}
Exact support multiplication gives precisely four row equations:
\begin{equation}\label{hodge:four-rows}
 \begin{aligned}
 \lambda c_{11}+\mu c_1&=0,&
 \lambda c_{29}+\mu c_{19}&=0,\\
 \lambda c_{51}+\mu c_{41}&=0,&
 \lambda c_{61}+\mu c_{51}&=0.
 \end{aligned}
\end{equation}
At the end of this round, no nonzero rational class satisfying the necessary
coefficient relations, no all-orders deformation, and no non-algebraicity
obstruction had been constructed. The next round tested this very direction
rather than treating the free complex coefficient system as a solution.
\src{C2}\src{C3}

\paragraph{Round III: the second-order claim and its exact scope.}
The third pass allowed arbitrary nonlinear correction at second order:
\begin{equation}\label{hodge:second-order-family}
 F_s=F_0+s h+s^2k+O(s^3),\qquad k\in S_{70},
\end{equation}
with $h$ as in \eqref{hodge:beta-pair}. Its recorded conclusion is that for
any nonzero rational $\alpha\in W$, no nonzero such $h$ can preserve the
Hodge condition through order two, regardless of $k$. This would in
particular exclude an analytic Hodge-preserving family with that initial
direction. The argument below is the archived mathematical argument,
supported by exact finite checks; it is not a claim about every possible
initial direction and not a Lean formalization.\src{C3}

\paragraph{Round III: period coordinates, nonvanishing, and reality.}
Let $\Omega$ denote the standard projective volume form and define
\begin{equation}\label{hodge:period-coordinates}
 \omega_t=\Res\left(\frac{x^{b(t)-\mathbf1}\Omega}{F_0^3}\right),
 \qquad G_t=\prod_{i=0}^{5}\Gamma\left(\frac{b_i(t)}{70}\right),
 \qquad K_t=\frac{G_t}{G_{70-t}}>0,
\end{equation}
where indices are taken modulo $70$. Use period coordinates
\[
 c_t=\int_{\PD(\alpha)}\omega_{70-t}.
\]
The source notes identify these with the earlier coefficient convention up
to a common nonzero normalization; alternatively the calculation can be
read directly in these period coordinates.

Every $c_t$ is nonzero for nonzero rational $\alpha\in W$. Galois conjugacy
forces every character projection of $\alpha$ to be nonzero, and cup
product pairs each eigenline nondegenerately with the inverse eigenline.
The explicit Fermat period formula on a spanning set of integral vanishing
cycles, attributed to Proposition 3.1 of Duque Franco--Villaflor Loyola,
gives
\begin{equation}\label{hodge:conjugation}
 \overline{\omega_t}=-K_t\omega_{70-t},\qquad
 c_{70-t}=-K_t\overline{c_t}.
\end{equation}
The sign comes from the common imaginary normalization constant; the
cyclotomic factors are conjugated by $t\mapsto70-t$. Rationality makes
$\PD(\alpha)$ real. The full normalization used in this formula is written
out in Round IV.\src{E9}\src{C3}

If a nonzero $h$ satisfies \eqref{hodge:four-rows}, then
$\lambda\mu\ne0$. Put $\rho=-\mu/\lambda$. In particular,
\[
 c_{29}=\rho c_{19},\qquad
 c_{51}=\rho c_{41},\qquad
 c_{61}=\rho c_{51}.
\]
Combining these identities with \eqref{hodge:conjugation} for the inverse
pairs $(19,51)$, $(29,41)$, and $(9,61)$ gives
\begin{equation}\label{hodge:rho}
 |\rho|^2=\frac{K_{19}}{K_{29}},\qquad
 \rho c_9=\frac{K_{19}^2}{K_9K_{29}}c_{19}.
\end{equation}
Only reality and nonvanishing of the character components are used at this
stage; rationality supplied both.

\paragraph{Round III: the exact Gamma ratio is four, not one.}
The crucial identity is
\begin{equation}\label{hodge:gamma-four}
 \frac{K_{19}^2}{K_9K_{29}}
 =\frac{G_{19}^{\,2}G_{61}G_{41}}
        {G_{51}^{\,2}G_9G_{29}}
 =4.
\end{equation}
It is certified by multiplication identities, not by a numerical Gamma
approximation. To keep two later uses of $M$ distinct, write
\[
 R_\Gamma(z)=\frac{\Gamma(z)}{\Gamma(1-z)},\qquad
 M_R(p,j)=
 \frac{\prod_{k=0}^{p-1}R_\Gamma((j+k70/p)/70)}
      {R_\Gamma(pj/70)}.
\]
Gauss's formula applied to numerator and denominator yields
\[
 M_R(p,j)=p^{1-2pj/70}.
\]
The source certificate is
\begin{equation}\label{hodge:gamma-four-certificate}
 \begin{aligned}
 \frac{K_{19}^2}{K_9K_{29}}
 &=\frac{M_R(2,1)M_R(2,3)M_R(2,4)M_R(2,8)M_R(2,10)M_R(2,17)}
       {M_R(2,5)M_R(2,9)M_R(2,13)M_R(2,16)}\\
 &\hspace{1cm}\cdot\frac{M_R(5,2)^2}{M_R(5,1)M_R(5,3)}.
 \end{aligned}
\end{equation}
The sum of the $j$ indices is $43$ on both sides of the $p=2$ quotient,
but there are six factors upstairs and four downstairs. Hence its value is
$2^2$. The exponent in the $p=5$ quotient is zero. The verifier compares the
individual $R_\Gamma$ exponents as integers and all prime exponents as exact
fractions. Thus
\begin{equation}\label{hodge:rho-four}
 \rho c_9=4c_{19}.
\end{equation}
The general Gamma identities are from NIST DLMF; the finite combination is
the specialized certificate from this pass.\src{E10}\src{C3}

\paragraph{Round III: a quadratic period immune to every degree-$70$ correction.}
Take
\[
 r=(35,28,2,12,5,58),\qquad \sum_i r_i=140,
\]
and the $F^3$ test form
\[
 \omega_r(s)=\Res\left(\frac{x^{r-\mathbf1}\Omega}{F_s^2}\right),
 \qquad x^{r-\mathbf1}=x_0^{34}x_1^{27}x_2x_3^{11}x_4^4x_5^{57}.
\]
The expansion is
\begin{equation}\label{hodge:quadratic-expansion}
 F_s^{-2}=F_0^{-2}-2shF_0^{-3}
 +s^2\bigl(3h^2F_0^{-4}-2kF_0^{-3}\bigr)+O(s^3).
\end{equation}
For a monomial $x^\gamma$ of $k$ to contribute to a block character at pole
three without further reduction, it would be necessary that
\[
 r+\gamma=b(t),\qquad b(t)\ge r\text{ coordinatewise}.
\]
The exact comparison of all $24$ characters finds none. Terms requiring
reduction from pole three have class in $F^3$ and also pair to zero with
$\alpha$. Thus the entire $k$ contribution vanishes, not just a finite
selection of correction monomials.

The three possible surviving terms of $h^2$ satisfy
\begin{align*}
 r+2\beta_2&=b(61)+70\mathbf e_5,\\
 r+\beta_1+\beta_2&=b(51)+70\mathbf e_5,\\
 r+2\beta_1&=b(41)+70\mathbf e_5.
\end{align*}
One Griffiths--Dwork reduction from pole four to pole three, including the
factor $3$ in \eqref{hodge:quadratic-expansion} and the mixed-term factor
$2$, gives the exact coefficient
\begin{equation}\label{hodge:quadratic-Q}
 Q=\frac{\mu^2c_9+12\lambda\mu c_{19}+11\lambda^2c_{29}}{35}.
\end{equation}
The three factors are respectively
$3(2/(70\cdot3))=1/35$,
$3\cdot2(12/(70\cdot3))=12/35$, and
$3(22/(70\cdot3))=11/35$.
Substituting the first-order relations and \eqref{hodge:rho-four} yields
\begin{equation}\label{hodge:nonzero-quadratic}
 Q=\frac{\lambda^2\rho}{35}(\rho c_9-c_{19})
   =\frac{3}{35}\lambda^2\rho c_{19}\ne0.
\end{equation}
This contradicts Hodge-period vanishing. Arbitrary higher-order corrections
cannot remove a nonzero second-order coefficient. Equivalently, the proposed
lift would require the Gamma ratio to be $1$, while its exact value is $4$.
This excludes the specified direction, not algebraic representatives of the
original class.\src{E7}\src{C3}

\paragraph{Round III: all six quadratic rows retained in the saved certificate.}
The verifier records six rows with no available linear correction. For each
row $r$, the following expression is the coefficient from $h^2$, before
imposing the first-order relations. A term indexed by power $j$ in the JSON
means $\lambda^j\mu^{2-j}c_t$ times its listed rational factor. The fifth
row is the one used above; the other rows are retained as part of the
computational record.\src{C3}

\begin{longtable}{@{}p{0.35\textwidth}p{0.61\textwidth}@{}}
\toprule
$r$ & Quadratic coefficient\\
\midrule
\endhead
$(3,18,4,68,13,34)$ & $\frac{2}{5}\mu^2{}c_{41}+\frac{4}{5}\lambda\mu{}c_{51}+\frac{2}{5}\lambda^2{}c_{61}$\\
$(7,28,30,26,47,2)$ & $\frac{11}{35}\mu^2{}c_{37}+\frac{3}{35}\lambda\mu{}c_{47}+\frac{23}{70}\lambda^2{}c_{57}$\\
$(15,48,12,12,45,8)$ & $\frac{2}{35}\mu^2{}c_{29}+\frac{1}{35}\lambda\mu{}c_{39}$\\
$(27,8,20,26,7,52)$ & $\frac{6}{35}\mu^2{}c_{17}+\frac{6}{35}\lambda\mu{}c_{27}+\frac{8}{35}\lambda^2{}c_{37}$\\
$(35,28,2,12,5,58)$ & $\frac{1}{35}\mu^2{}c_{9}+\frac{12}{35}\lambda\mu{}c_{19}+\frac{11}{35}\lambda^2{}c_{29}$\\
$(41,8,6,54,21,10)$ & $\frac{1}{5}\mu^2{}c_{3}+\frac{2}{5}\lambda\mu{}c_{13}+\frac{1}{5}\lambda^2{}c_{23}$\\
\bottomrule
\end{longtable}

\paragraph{Round III: bounded extension to a common-monomial family.}
The chosen pair factors as
\[
 h=x^\eta\bigl(\lambda x_1^{10}x_4^{20}x_5^{10}
                   +\mu x_0^{10}x_2^{30}\bigr),
 \qquad\eta=(3,1,1,15,3,7),\quad|\eta|=30.
\]
Among all $\binom{35}{5}=324,632$ degree-$30$ common monomials, $170$ give
exactly the same four first-order equations. For $78$ of those $170$, the
same uncorrectable quadratic-row construction works. The saved C++ census
runs over all degree-$30$ compositions; the independent Python implementation
first restricts to the necessary box
\[
 \eta\le(8,19,3,27,8,7).
\]
The counts agree. This does not exclude the remaining common monomials, all
two-monomial tangents, or other varieties. The full $78$ witness pairs are
preserved in the original JSON and in the certificate supplement later in
this document.\src{C3}

\paragraph{Round IV: a distinguished nonzero rational Hodge class.}
The final counterexample pass changed from deformations to normalized
periods. Retain \eqref{hodge:period-coordinates} and use the normalization of
Proposition 3.1 in the cited Fermat period paper. For a standard vanishing
cycle $\delta_v$, with $v_0=0$, the formula used in the archive is
\begin{equation}\label{hodge:full-period-formula}
 \int_{\delta_v}\omega_t
 =\frac{G_t}{2\cdot70^5(2\pi i)}
   \prod_{i=0}^{5}
   \bigl(\zeta^{b_i(t)(v_i+1)}-\zeta^{b_i(t)v_i}\bigr).
\end{equation}
Those cycles span primitive rational homology. Define the ordinary,
untwisted Betti class
\begin{equation}\label{hodge:explicit-alpha}
 \boxed{\quad
 \alpha=70^4\sum_{t\in U}\frac{2\pi i}{G_t}\,\omega_t.
 \quad}
\end{equation}
Its period on $\delta_v$ is $1/140$ times the trace from $K$ to $\Q$ of the
cyclotomic product in \eqref{hodge:full-period-formula}, hence rational.
The residue classes are primitive and pair to zero with the ambient class,
so the spanning-period test gives rationality. Each summand is of type
$(2,2)$, so $\alpha$ is Hodge.\src{E9}\src{C4}

For $v=(0,\ldots,0)$, the saved exact calculation is
\begin{equation}\label{hodge:trace-normalization}
 \Tr_{K/\Q}\prod_{i=0}^{5}(\zeta^{a_i}-1)=140,
 \qquad \int_{\delta_0}\alpha=1.
\end{equation}
The verifier evaluates this trace by Ramanujan sums rather than by
floating-point Gamma values. More explicitly, with
\[
 c_{70}(k)=\sum_{t\in U}\zeta^{tk}
 =\mu\left(\frac{70}{\gcd(k,70)}\right)
   \frac{\varphi(70)}{\varphi(70/\gcd(k,70))},
\]
it checks the finite integer identity
\[
 \sum_{\varepsilon\in\{0,1\}^6}
 (-1)^{6-|\varepsilon|}
 c_{70}\left(\sum_i\varepsilon_i a_i\right)=140.
\]
Here $\mu$ is the M\"obius function and $\varphi$ Euler's totient, as
implemented in the supplied verifier. Thus the recorded class is explicitly
nonzero:
\[
 0\ne\alpha\in H^4(X,\Q)\cap H^{2,2}(X).
\]
Rationality and nonvanishing are not non-algebraicity.

\paragraph{Round IV: the proposed transcendence obstruction and its failure.}
Because $X$ is defined over $\Q$, the final pass used the necessary condition
that an algebraic class can be represented after spreading and specializing
its cycle to an algebraic point of a suitable Hilbert or Chow component.
Its Tate-normalized algebraic de Rham expression then has algebraic
coefficients. This is the intended necessary-condition argument, not a
sufficient criterion for algebraicity. The archive distinguishes the
untwisted Betti normalization from the Tate-normalized de Rham one.\src{C4}

Set
\[
 C_t=\frac{G_t}{(2\pi i)^3}.
\]
Then the displayed normalization gives
\begin{equation}\label{hodge:tate-normalization}
 (2\pi i)^2\alpha
 =70^4\sum_{t\in U}C_t^{-1}\omega_t.
\end{equation}
A transcendental coefficient in the indicated algebraic de Rham basis would
have supplied a non-algebraicity obstruction. But all the $C_t$ here are
algebraic. The general algebraicity of these normalized Fermat Gamma
products is attributed to Deligne's established theory; the new record is
an explicit specialization, not a new general theorem.\src{E11}\src{C4}

\paragraph{Round IV: a short exact Gamma certificate for the selected character.}
Write $G=G_1$, $G_-=G_{69}$. The recorded identity is
\begin{equation}\label{hodge:gamma-radical}
 \boxed{\quad \frac{G}{G_-}
  =2^{16/35}5^{3/14}7^{2/5}.\quad}
\end{equation}
Positive real powers are intended. Define a different multiplication ratio
from that of Round III:
\[
 M_\Gamma(p,j)=
 \frac{\prod_{k=0}^{p-1}\Gamma((j+k70/p)/70)}{\Gamma(pj/70)},
 \qquad p\in\{2,5,7\},\quad1\le j<70/p.
\]
Gauss multiplication gives
\[
 M_\Gamma(p,j)=(2\pi)^{(p-1)/2}p^{1/2-pj/70}.
\]
The exact cancellation certificate is
\begin{equation}\label{hodge:short-gamma-certificate}
 \frac{G}{G_-}
 =\frac{M_\Gamma(2,1)M_\Gamma(2,26)M_\Gamma(7,2)M_\Gamma(7,4)}
 {M_\Gamma(2,9)M_\Gamma(2,34)
  M_\Gamma(5,4)M_\Gamma(5,8)M_\Gamma(5,12)}.
\end{equation}
The power of $2\pi$ cancels, leaving the exact prime exponents in
\eqref{hodge:gamma-radical}. The verifier checks every Gamma exponent as an
integer and every resulting prime exponent as a rational number.\src{E10}\src{C4}

The saved file \texttt{all\_gamma\_certificates.json} contains a multiplication
certificate for each of the twelve inverse pairs. Its full exponent table
is reproduced next: if a row is $(t,e_2,e_5,e_7)$, then
$G_t/G_{70-t}=2^{e_2}5^{e_5}7^{e_7}$. The other twelve rows are obtained by
negating the exponents. Reflection then gives algebraicity of every $C_t$.

\begin{longtable}{rrrr}
\toprule
$t$ & $e_2$ & $e_5$ & $e_7$\\
\midrule
\endhead
$1$ & $\frac{16}{35}$ & $\frac{3}{14}$ & $\frac{2}{5}$\\
$3$ & $\frac{48}{35}$ & $\frac{9}{14}$ & $\frac{1}{5}$\\
$9$ & $\frac{4}{35}$ & $-\frac{1}{14}$ & $-\frac{2}{5}$\\
$11$ & $\frac{36}{35}$ & $-\frac{9}{14}$ & $\frac{2}{5}$\\
$13$ & $-\frac{2}{35}$ & $-\frac{3}{14}$ & $\frac{1}{5}$\\
$17$ & $-\frac{8}{35}$ & $\frac{9}{14}$ & $-\frac{1}{5}$\\
$19$ & $\frac{24}{35}$ & $\frac{1}{14}$ & $-\frac{2}{5}$\\
$23$ & $\frac{18}{35}$ & $-\frac{1}{14}$ & $\frac{1}{5}$\\
$27$ & $\frac{12}{35}$ & $-\frac{3}{14}$ & $-\frac{1}{5}$\\
$29$ & $-\frac{26}{35}$ & $\frac{3}{14}$ & $-\frac{2}{5}$\\
$31$ & $\frac{6}{35}$ & $\frac{9}{14}$ & $\frac{2}{5}$\\
$33$ & $-\frac{32}{35}$ & $\frac{1}{14}$ & $\frac{1}{5}$\\
\bottomrule
\end{longtable}

These certificates independently recover
\[
 2e_p(19)-e_p(9)-e_p(29)
 =\begin{cases}2,&p=2,\\0,&p=5,7,\end{cases}
\]
and hence \eqref{hodge:gamma-four}. The general multiplication and reflection
formulas are classical; the stated finite combinations and their output
belong to the recorded investigation.

\paragraph{Round IV: the explicit degree-$70$ Kummer equation.}
Put
\begin{equation}\label{hodge:kummer-definitions}
 \begin{gathered}
 K=\Q(\zeta),\qquad
 P=\prod_{i=0}^{5}(1-\zeta^{a_i}),\qquad C=\frac{G}{(2\pi i)^3},\\
 D=-\frac{2^{16}5^7\sqrt5\,7^{14}}{P^{35}}.
 \end{gathered}
\end{equation}
Here $D$ is the radical's right-hand side, not the earlier tangent ideal
$\mathfrak d$ or the determinant in the final witness search. The positive
square root is in $K$, with explicit formula
\[
 \sqrt5=1+2(\zeta^{14}+\zeta^{56}).
\]
The verifier checks its square is $5$ in $\Z[\zeta]$.
Since $\sum_i a_i=210$,
\[
 P=64\prod_i\sin\left(\frac{\pi a_i}{70}\right)>0.
\]
Euler reflection gives
\[
 GG_-=\frac{\pi^6}{\prod_i\sin(\pi a_i/70)},\qquad
 C^2=-\frac{G/G_-}{P}.
\]
Raising to the $35$th power and applying \eqref{hodge:gamma-radical} yields
\begin{equation}\label{hodge:kummer}
 \boxed{\quad C^{70}
 =-\frac{2^{16}5^7\sqrt5\,7^{14}}{P^{35}}
 =D\in K.\quad}
\end{equation}
This supplies the upper bound $[K(C):K]\le70$.\src{E10}\src{C4}

\paragraph{Round IV: the valuation proof of exact degree.}
The inherited argument for equality of the degree is included in full.
Since $K=\Q(\zeta_{35})$, the prime $2$ is unramified in $K$ and the
ramification index at $5$ is $4$. Use integer-normalized valuations on $K$.
The factors $1-\zeta^{a_i}$ in $P$ have respective root orders
\[
 70,\ 7,\ 35,\ 5,\ 70,\ 35.
\]
A factor $1-\xi$ for a root of unity of non-prime-power order is a unit;
at primes above $5$, the order-$5$ factor has valuation $1$ and the
order-$7$ factor is a unit. Consequently
\[
 v_{\mathfrak p}(P)=0\quad(\mathfrak p\mid2),\qquad
 v_{\mathfrak q}(P)=1\quad(\mathfrak q\mid5),
\]
and \eqref{hodge:kummer-definitions} gives
\begin{equation}\label{hodge:valuations}
 v_{\mathfrak p}(D)=16,\qquad
 v_{\mathfrak q}(D)=4\cdot7+2-35=-5.
\end{equation}
A $70$th root has valuations $8/35$ and $-1/14$ on extending these
valuations. A field containing it must have ramification indices divisible
by $35$ and $14$, respectively. Its degree is divisible by
$\operatorname{lcm}(35,14)=70$. Together with the upper bound this yields
\begin{equation}\label{hodge:exact-kummer-degree}
 [K(C):K]=70.
\end{equation}
Since $K$ contains all $70$th roots of unity, adjoining one root splits
$T^{70}-D$ and the extension is cyclic. This proof concerns a field generated
by a period factor, not the existence or nonexistence of an algebraic
surface. The supplied program verifies the arithmetic least-common-multiple
step; that check alone is not a proof of the valuation assertions preceding
it.\src{C4}

\paragraph{Round IV: independent Jacobi--Kummer comparisons at thirty-two primes.}
For each tested $p\equiv1\pmod{70}$, take the indicated primitive root $g$
and the embedding $\zeta\mapsto z=g^{(p-1)/70}$. Reducing $D$ modulo $p$
gives
\[
 u=D^{(p-1)/70}\pmod p=z^e.
\]
The saved verifier separately computes the Jacobi eigenvalue in
$\Z[\zeta]$ and checks $\lambda_p=p^2\zeta^e$. All $32$ split primes
from $71$ through $6301$ agree. The four earlier primes are included, not
replaced by newly selected favorable primes. A finite list of matching
primes is not a proof of equality of global characters; the conversation
invokes Deligne's Theorem 7.15 for the general period--Galois relationship.
Inverting a Frobenius convention leaves the character kernel unchanged.
\src{E11}\src{C4}

The full saved comparison table follows. All entries are finite-field or
integer computations; $D_p$ denotes $D\bmod p$.

\begingroup\small

\begin{longtable}{rrrrrr}
\toprule
$p$ & $g$ & $z$ & $D_p$ & $e$ & Order\\
\midrule
\endhead
$71$ & $7$ & $7$ & $20$ & $40$ & $7$\\
$211$ & $2$ & $8$ & $87$ & $12$ & $35$\\
$281$ & $3$ & $81$ & $149$ & $32$ & $35$\\
$421$ & $2$ & $64$ & $109$ & $50$ & $7$\\
$491$ & $2$ & $128$ & $19$ & $57$ & $70$\\
$631$ & $3$ & $122$ & $166$ & $64$ & $35$\\
$701$ & $2$ & $323$ & $546$ & $34$ & $35$\\
$911$ & $17$ & $66$ & $128$ & $0$ & $1$\\
$1051$ & $7$ & $322$ & $895$ & $21$ & $10$\\
$1471$ & $6$ & $526$ & $1040$ & $3$ & $70$\\
$2311$ & $3$ & $2050$ & $1323$ & $59$ & $70$\\
$2381$ & $3$ & $934$ & $223$ & $55$ & $14$\\
$2521$ & $17$ & $1336$ & $801$ & $17$ & $70$\\
$2591$ & $7$ & $2365$ & $457$ & $18$ & $35$\\
$2731$ & $3$ & $1093$ & $2352$ & $37$ & $70$\\
$2801$ & $3$ & $1481$ & $2761$ & $58$ & $35$\\
$3011$ & $2$ & $1686$ & $2472$ & $11$ & $70$\\
$3221$ & $10$ & $2879$ & $133$ & $13$ & $70$\\
$3361$ & $22$ & $288$ & $1230$ & $10$ & $7$\\
$3571$ & $2$ & $3413$ & $3076$ & $13$ & $70$\\
$3851$ & $2$ & $83$ & $551$ & $62$ & $35$\\
$4201$ & $11$ & $98$ & $526$ & $44$ & $35$\\
$4271$ & $7$ & $2359$ & $1240$ & $12$ & $35$\\
$4481$ & $3$ & $635$ & $570$ & $35$ & $2$\\
$4621$ & $2$ & $4030$ & $2656$ & $0$ & $1$\\
$4691$ & $2$ & $3297$ & $502$ & $38$ & $35$\\
$4831$ & $3$ & $4498$ & $4497$ & $35$ & $2$\\
$5531$ & $10$ & $815$ & $2566$ & $11$ & $70$\\
$5741$ & $2$ & $3891$ & $5667$ & $60$ & $7$\\
$5881$ & $31$ & $5741$ & $2192$ & $19$ & $70$\\
$6091$ & $7$ & $466$ & $2962$ & $58$ & $35$\\
$6301$ & $10$ & $1108$ & $2690$ & $36$ & $35$\\
\bottomrule
\end{longtable}

\endgroup

The archive also records the reductions of $\sqrt5$ and $P$ at every prime,
and those full machine-readable rows remain preserved in the embedded
original archive. The explicit radical accounts for the observed finite
orders $7$, $35$, and $70$. It does not produce an eigenvalue of infinite
order or an obstruction to cycles over $\C$.

\paragraph{Round IV: enlarging the formal character calculations.}
The final pass enumerated nondecreasing quadruples with entries in
$\{1,\ldots,69\}$ and sum $140$, and found exactly $700$ satisfying
\[
 \sum_{i=1}^{4}\langle tq_i\rangle_{70}=140
 \quad\text{for every }t\in U.
\]
Their multiplicity vectors were adjoined to the pair-and-standard generator
list. The chosen sextuple's multiplicity vector still lies outside the
enlarged lattice: reduction modulo $2$ is outside the generator span. The
reported mod-$2$ rank is $55$. This rank is over $\F_2$, not the
characteristic-zero lattice rank $57$ reported earlier. The quadruples
represent divisor-type Hodge blocks on Fermat surfaces; their formal
multiplicity vectors are not divisor classes on the original fourfold.
\src{C4}

For each $1\le f\le100$, the code then inflated the character to
\[
 m=70f,\qquad a_f=fa,
\]
corresponding to the finite power map
\[
 X_m\longrightarrow X_{70},\qquad
 [x_0:\cdots:x_5]\longmapsto[x_0^f:\cdots:x_5^f].
\]
At all $100$ tested levels, through degree $7000$, $\nu(a_f)$ remains
outside the specified pair-and-standard lattice, again certified by exact
mod-$2$ linear algebra. The full rank record is included later. This is not
a result about all covers, all constructions, or the full cycle-class image.
The two-torsion phenomenon in formal multiplicity vectors must not be
confused with an integral cycle-class obstruction.\src{E2}\src{E3}\src{C4}

\paragraph{The all-or-nothing question for this rational block.}
Let $e_W$ be the rational character projector onto $W$. The coordinate-scaling
automorphisms are algebraic, and the projector is a rational combination of
their graph correspondences. Thus
\[
 e_W\bigl(\cl(\CH^2(X_{\C})_{\Q})\bigr)
\]
is an automorphism-stable rational subspace of $W$ consisting of algebraic
cycle classes. The rational character representation on $W$ is irreducible;
therefore this image is either $0$ or all of $W$. A single nonzero projected
algebraic class would generate the whole block under the algebraic
automorphisms, including \eqref{hodge:explicit-alpha}.\src{C4}\src{C5}

For the specified nonzero rational class to be a counterexample, the
unresolved assertion is precisely
\begin{equation}\label{hodge:remaining-alpha}
 \alpha\notin\cl\bigl(\CH^2(X_{\C})_{\Q}\bigr),
\end{equation}
or, equivalently for this block,
\begin{equation}\label{hodge:remaining-block}
 \boxed{\quad e_W\bigl(\cl(\CH^2(X_{\C})_{\Q})\bigr)=0.\quad}
\end{equation}
No round established this equality or its opposite. Failure of selected
construction tests, finite field-of-definition restrictions, failure of a
deformation to remain Hodge, and algebraicity of normalized periods do not
settle \eqref{hodge:remaining-block}.

\paragraph{Reported breakthroughs: retained as provenance, not mathematical premises.}
After the deformation rounds, the user reported a rumor that Claude had
found a counterexample and was still Lean-formalizing it. The subsequent
assistant response said it had not located a public construction or
Lean-checked proof, and used an ``as of 14 September 2026'' marker. It also
mentioned a purported 12 September status discussion by Claire Voisin and
an institutional page said to remain marked unsolved. Those are assertions
made in the earlier conversation; no underlying manuscript, formal proof,
or text of the purported Voisin discussion is among the supplied research
artifacts. They are not independently verified in this export.\src{C5}

The later user message supplied the following wording from the linked
MTSlive post:
\begin{quote}
``SITUATION BREWING: OpenAI is nearing a solution to the Hodge Conjecture,
one of the remaining Millennium problems, and expects the proof to be
finished soon, per The Information.''
\end{quote}
The supplied post is retained in the source list as \src{N1}. The previous
assistant then said it had found a matching subscriber-only article and an
indexed public statement attributed to Stephanie Palazzolo, and described
that as more identifiable reporting than the earlier anonymous rumor. Its
response used a different ``as of 17 September 2026'' marker. It did not
provide a mathematical theorem statement, a Hodge manuscript, a reproducible
formalization, or the article's full text. It also mentioned a separate
Navier--Stokes item on OpenAI's research index. That separate claim is
retained here only as part of what was previously said, not as an endorsed
or newly checked announcement.\src{N2}\src{N3}\src{C5}

The inherited 14 and 17 September markers are not export-time verification
dates. This document is dated 27 September 2026 and makes no fresh claim
about the current public status of any purported solution. The available
conversation does not supply an actual Hodge proof or counterexample from
Claude, OpenAI, the reporter, or the social-media account. Unpublished work
is neither established nor excluded by that absence.

The mathematical distinction made in the conversation remains important.
A ``solution'' could be a positive proof or a counterexample. A positive
proof of \eqref{hodge:statement} would make every genuine rational Hodge
class in $W$ algebraic and rule out this attempted counterexample. A
counterexample somewhere else would not itself prove non-algebraicity of
our particular $\alpha$. These are logical consequences of the statement's
quantifiers, not inferences about the content of any unpublished work.
The conversation also expressly acknowledged no access to OpenAI's
unpublished research and no evidence connecting these calculations to the
reported work.\src{C5}

The earlier replies offered alerts or daily checks, but no configured
monitoring task or subsequent alert is present in the supplied record.
These conversational offers are not mathematical progress or evidence of
a public result.

\paragraph{The user's existence heuristic and the change of search direction.}
The user observed that even knowing a solution exists may help find it,
however difficult it is. The response distinguished using possible
existence as a search heuristic from treating it as an unproved mathematical
premise. It then reversed the local task: rather than accumulating failures
of familiar constructions, seek the missing surface whose class would make
this candidate algebraic. This does not assume that a reported solution
exists or that the full conjecture is true. It asks what a successful
positive witness for this one block would look like.\src{C5}

By the all-or-nothing argument, a sufficient target is
\begin{equation}\label{hodge:positive-witness}
 \boxed{\quad
 Z\subset X\text{ an algebraic surface with }e_W\cl(Z)\ne0.
 \quad}
\end{equation}
Finding it would settle algebraicity of the entire selected block and
eliminate this block as a counterexample. It would not prove the Hodge
conjecture for other blocks or other varieties.

\paragraph{Final strategy: a complete-intersection surface certificate.}
The last mathematical response invokes the complete-intersection cycle-class
formula attributed to Villaflor Loyola. It seeks homogeneous polynomials
$f_i,g_i$ satisfying
\begin{equation}\label{hodge:witness-decomposition}
 \boxed{\qquad F_0=f_1g_1+f_2g_2+f_3g_3.\qquad}
\end{equation}
Let $d_i=\deg f_i$, so $\deg g_i=70-d_i$. One must also require that
\[
 Z=V(f_1,f_2,f_3)\subset\mathbb P^5
\]
be a complete-intersection surface, with the usual cycle multiplicities
allowed. Merely writing an identity in polynomials does not certify the
required dimension or regular-sequence condition. Identity
\eqref{hodge:witness-decomposition} then puts $Z$ inside $X$.

Define the $6\times6$ Jacobian determinant
\begin{equation}\label{hodge:witness-determinant}
 D_Z=\det\left(
 \frac{\partial(f_1,g_1,f_2,g_2,f_3,g_3)}
      {\partial(x_0,x_1,x_2,x_3,x_4,x_5)}\right).
\end{equation}
Its degree is
\[
 \sum_{i=1}^{3}\bigl((d_i-1)+(70-d_i-1)\bigr)=204.
\]
The cited formula identifies the primitive cycle class with a nonzero scalar
multiple of the residue represented by $D_Z$ modulo the Jacobian ideal.
This is the named theorem used as the bridge from a polynomial construction
to a cycle class; the specialized witness target below is from the last
response.\src{E12}\src{C5}

Since $J=(x_0^{69},\ldots,x_5^{69})$, the proposed sufficient coefficient
condition is
\begin{equation}\label{hodge:witness-coefficient}
 \boxed{\quad
 [x_0^0x_1^{19}x_2^{23}x_3^{41}x_4^{60}x_5^{61}]\,D_Z\ne0.
 \quad}
\end{equation}
The bracket denotes the coefficient of that monomial. Its exponent vector
is $a-\mathbf1$, its degree is $204$, and each exponent is at most $68$,
so the monomial survives in $R=S/J$. Under the Fermat character
identification, a nonzero coefficient supplies the required nonzero block
component. The search therefore has an explicit positive witness target:
the polynomial identity, the complete-intersection condition, and one
nonzero determinant coefficient.\src{E2}\src{E12}\src{C5}

\paragraph{Final strategy: why the separated-pair construction cannot work.}
The first family tested factors the six variables into three pairs and
expresses each paired sum of two $70$th powers as one product $f_ig_i$.
For such a separated-pair construction, the determinant contribution for
each pair has degree $68$ in just those two variables. A target monomial
can occur only if its corresponding character entries obey
\[
 (a_i-1)+(a_j-1)=68,\qquad\text{that is, }a_i+a_j=70
\]
within each pair. The selected character has no complementary pair.
The last response records checking all $15$ coordinate pairings and all
$24$ conjugate characters, with no surviving separated-pair witness.
There is no separate code archive for this final chat-only check.
It excludes that family, not polynomials mixing variables across pairs.
\src{C5}

\paragraph{Final strategy: finite degree profiles and the unsolved coefficient system.}
Swapping $f_i$ with $g_i$ leaves
\eqref{hodge:witness-decomposition} unchanged and changes only the sign of
$D_Z$. Reordering the three pairs therefore allows positive degree profiles
\[
 1\le d_1\le d_2\le d_3\le35.
\]
There are
\begin{equation}\label{hodge:degree-profiles}
 \binom{37}{3}=7,770
\end{equation}
such profiles. For each, comparison of coefficients in
\eqref{hodge:witness-decomposition} yields polynomial equations in the
coefficients of the six unknown polynomials. If $c$ denotes the coefficient
in \eqref{hodge:witness-coefficient}, an auxiliary variable $u$ can encode
nonvanishing by $uc=1$. The geometric complete-intersection requirement
must still be checked; it is not silently supplied by that auxiliary
equation.

This is a finite algebraic feasibility formulation for the restricted
witness class, but unrestricted polynomials can involve millions of
coefficients. The conversation did not solve these systems, did not exhaust
all $7,770$ profiles, and did not find six polynomials meeting the stated
conditions. Even excluding every complete-intersection witness would not
exclude arbitrary algebraic surfaces or rational combinations of cycles.
The final contribution is a concrete witness target, not a successful
construction or a proof of the universal conjecture.\src{C5}

\paragraph{Consolidated logical audit.}
The following table summarizes the progression without conflating a
necessary condition with a sufficient one. All entries concern the chosen
candidate unless explicitly stated otherwise.
\begin{longtable}{@{}L{0.24\textwidth}L{0.36\textwidth}L{0.34\textwidth}@{}}
\toprule
Approach & Recorded result & What remains unproved\\
\midrule\endhead
Fermat characters & A $24$-dimensional rational $(2,2)$ block; later an explicit nonzero rational class & Whether any nonzero class in the block is algebraic\\
Elementary and formal construction tests & Residual cases; selected multiplicity vector outside specified lattices & Exclusion of all algebraic cycle constructions\\
Frobenius and Hecke characters & Finite-order values, exact order $70$, finite field-of-definition restriction & Non-algebraicity over $\C$\\
Straight monomial deformations & Every tested basis direction escapes by order eight & Exclusion of arbitrary curved or multivariate Hodge-preserving deformations\\
Whole-block nonlinear family & Inherited local power-family classification and correspondence invariance & Algebraicity or non-algebraicity of the fixed pulled-back block\\
Specified two-monomial tangent & Uncorrectable second-order coefficient $3\lambda^2\rho c_{19}/35$ & Exclusion of all other individual-class directions\\
Normalized periods & Exact Gamma identities and cyclic degree-$70$ radical field & Construction of an algebraic surface, or a transcendence obstruction\\
One hundred inflated levels & A formal mod-$2$ obstruction at each specified level & A statement about every finite cover or the full cycle-class image\\
Complete intersections & Polynomial identity and determinant-coefficient witness target; paired family fails & A successful six-polynomial witness or a complete search\\
Reported breakthrough & User-supplied report and inherited discussion, no supplied proof & The actual statement, argument, and correctness of any unpublished work\\
\bottomrule
\end{longtable}

\paragraph{Certificate supplement: higher-order monomial directions.}
The continuation archive includes the complete CSV of directions whose first
escape order is at least five. The following table retains every row in that
file. The earlier census already records the full counts for orders one
through four.\src{C2}

Here $j$ is the first escape order, $t$ the witness unit, $q$ the witness pole, and support the number of nonzero coordinates of $\beta$.

\begingroup\small\setlength{\tabcolsep}{4pt}

\begin{longtable}{crrrrc}
\toprule
$\beta$ & Support & $j$ & $t$ & $q$ & $r$\\
\midrule
\endhead
$(0,0,0,67,0,3)$ & $2$ & $5$ & $3$ & $2$ & $(3,60,2,1,43,31)$\\
$(0,0,0,67,1,2)$ & $3$ & $5$ & $3$ & $2$ & $(3,60,2,1,38,36)$\\
$(0,0,0,67,2,1)$ & $3$ & $5$ & $3$ & $2$ & $(3,60,2,1,33,41)$\\
$(0,0,0,67,3,0)$ & $2$ & $5$ & $3$ & $2$ & $(3,60,2,1,28,46)$\\
$(0,0,0,68,0,2)$ & $2$ & $8$ & $3$ & $2$ & $(3,60,2,2,43,30)$\\
$(0,0,0,68,1,1)$ & $3$ & $8$ & $3$ & $2$ & $(3,60,2,2,35,38)$\\
$(0,0,0,68,2,0)$ & $2$ & $8$ & $3$ & $2$ & $(3,60,2,2,27,46)$\\
$(0,0,1,67,0,2)$ & $3$ & $5$ & $13$ & $2$ & $(13,50,27,1,23,26)$\\
$(0,0,1,67,1,1)$ & $4$ & $5$ & $13$ & $2$ & $(13,50,27,1,18,31)$\\
$(0,0,1,67,2,0)$ & $3$ & $5$ & $13$ & $2$ & $(13,50,27,1,13,36)$\\
$(0,0,1,68,0,1)$ & $3$ & $8$ & $13$ & $2$ & $(13,50,24,2,23,28)$\\
$(0,0,1,68,1,0)$ & $3$ & $8$ & $13$ & $2$ & $(13,50,24,2,15,36)$\\
$(0,0,2,67,0,1)$ & $3$ & $5$ & $13$ & $2$ & $(13,50,22,1,23,31)$\\
$(0,0,2,67,1,0)$ & $3$ & $5$ & $13$ & $2$ & $(13,50,22,1,18,36)$\\
$(0,0,2,68,0,0)$ & $2$ & $8$ & $13$ & $2$ & $(13,50,16,2,23,36)$\\
$(0,0,3,67,0,0)$ & $2$ & $5$ & $13$ & $2$ & $(13,50,17,1,23,36)$\\
$(0,1,0,67,0,2)$ & $3$ & $5$ & $3$ & $2$ & $(3,55,2,1,43,36)$\\
$(0,1,0,67,1,1)$ & $4$ & $5$ & $3$ & $2$ & $(3,55,2,1,38,41)$\\
$(0,1,0,67,2,0)$ & $3$ & $5$ & $3$ & $2$ & $(3,55,2,1,33,46)$\\
$(0,1,0,68,0,1)$ & $3$ & $8$ & $3$ & $2$ & $(3,52,2,2,43,38)$\\
$(0,1,0,68,1,0)$ & $3$ & $8$ & $3$ & $2$ & $(3,52,2,2,35,46)$\\
$(0,1,1,67,0,1)$ & $4$ & $5$ & $13$ & $2$ & $(13,45,27,1,23,31)$\\
$(0,1,1,67,1,0)$ & $4$ & $5$ & $13$ & $2$ & $(13,45,27,1,18,36)$\\
$(0,1,1,68,0,0)$ & $3$ & $8$ & $13$ & $2$ & $(13,42,24,2,23,36)$\\
$(0,1,2,67,0,0)$ & $3$ & $5$ & $13$ & $2$ & $(13,45,22,1,23,36)$\\
$(0,2,0,67,0,1)$ & $3$ & $5$ & $3$ & $2$ & $(3,50,2,1,43,41)$\\
$(0,2,0,67,1,0)$ & $3$ & $5$ & $3$ & $2$ & $(3,50,2,1,38,46)$\\
$(0,2,0,68,0,0)$ & $2$ & $8$ & $3$ & $2$ & $(3,44,2,2,43,46)$\\
$(0,2,1,67,0,0)$ & $3$ & $5$ & $13$ & $2$ & $(13,40,27,1,23,36)$\\
$(0,3,0,67,0,0)$ & $2$ & $5$ & $3$ & $2$ & $(3,45,2,1,43,46)$\\
$(0,68,0,0,0,2)$ & $2$ & $6$ & $3$ & $2$ & $(3,2,2,56,43,34)$\\
$(0,68,0,0,1,1)$ & $3$ & $6$ & $3$ & $2$ & $(3,2,2,56,37,40)$\\
$(0,68,0,0,2,0)$ & $2$ & $6$ & $3$ & $2$ & $(3,2,2,56,31,46)$\\
$(0,68,0,1,0,1)$ & $3$ & $6$ & $3$ & $2$ & $(3,2,2,50,43,40)$\\
$(0,68,0,1,1,0)$ & $3$ & $6$ & $3$ & $2$ & $(3,2,2,50,37,46)$\\
$(0,68,0,2,0,0)$ & $2$ & $6$ & $3$ & $2$ & $(3,2,2,44,43,46)$\\
$(0,68,1,0,0,1)$ & $3$ & $6$ & $31$ & $2$ & $(31,2,38,42,1,26)$\\
$(0,68,1,0,1,0)$ & $3$ & $6$ & $17$ & $2$ & $(17,2,52,14,51,4)$\\
$(0,68,1,1,0,0)$ & $3$ & $6$ & $17$ & $2$ & $(17,2,52,8,57,4)$\\
$(0,68,2,0,0,0)$ & $2$ & $6$ & $17$ & $2$ & $(17,2,46,14,57,4)$\\
$(1,0,0,67,0,2)$ & $3$ & $5$ & $13$ & $2$ & $(8,50,32,1,23,26)$\\
$(1,0,0,67,1,1)$ & $4$ & $5$ & $13$ & $2$ & $(8,50,32,1,18,31)$\\
$(1,0,0,67,2,0)$ & $3$ & $5$ & $13$ & $2$ & $(8,50,32,1,13,36)$\\
$(1,0,0,68,0,1)$ & $3$ & $8$ & $13$ & $2$ & $(5,50,32,2,23,28)$\\
$(1,0,0,68,1,0)$ & $3$ & $8$ & $13$ & $2$ & $(5,50,32,2,15,36)$\\
$(1,0,1,67,0,1)$ & $4$ & $5$ & $13$ & $2$ & $(8,50,27,1,23,31)$\\
$(1,0,1,67,1,0)$ & $4$ & $5$ & $13$ & $2$ & $(8,50,27,1,18,36)$\\
$(1,0,1,68,0,0)$ & $3$ & $8$ & $13$ & $2$ & $(5,50,24,2,23,36)$\\
$(1,0,2,67,0,0)$ & $3$ & $5$ & $13$ & $2$ & $(8,50,22,1,23,36)$\\
$(1,1,0,67,0,1)$ & $4$ & $5$ & $13$ & $2$ & $(8,45,32,1,23,31)$\\
$(1,1,0,67,1,0)$ & $4$ & $5$ & $13$ & $2$ & $(8,45,32,1,18,36)$\\
$(1,1,0,68,0,0)$ & $3$ & $8$ & $13$ & $2$ & $(5,42,32,2,23,36)$\\
$(1,1,1,67,0,0)$ & $4$ & $5$ & $13$ & $2$ & $(8,45,27,1,23,36)$\\
$(1,2,0,67,0,0)$ & $3$ & $5$ & $13$ & $2$ & $(8,40,32,1,23,36)$\\
$(1,68,0,0,0,1)$ & $3$ & $6$ & $31$ & $2$ & $(25,2,44,42,1,26)$\\
$(1,68,0,0,1,0)$ & $3$ & $6$ & $17$ & $2$ & $(11,2,58,14,51,4)$\\
$(1,68,0,1,0,0)$ & $3$ & $6$ & $17$ & $2$ & $(11,2,58,8,57,4)$\\
$(1,68,1,0,0,0)$ & $3$ & $6$ & $17$ & $2$ & $(11,2,52,14,57,4)$\\
$(2,0,0,67,0,1)$ & $3$ & $5$ & $13$ & $2$ & $(3,50,32,1,23,31)$\\
$(2,0,0,67,1,0)$ & $3$ & $5$ & $13$ & $2$ & $(3,50,32,1,18,36)$\\
$(2,0,0,68,0,0)$ & $2$ & $8$ & $23$ & $2$ & $(7,40,62,2,3,26)$\\
$(2,0,1,67,0,0)$ & $3$ & $5$ & $13$ & $2$ & $(3,50,27,1,23,36)$\\
$(2,1,0,67,0,0)$ & $3$ & $5$ & $13$ & $2$ & $(3,45,32,1,23,36)$\\
$(2,68,0,0,0,0)$ & $2$ & $6$ & $17$ & $2$ & $(5,2,58,14,57,4)$\\
$(3,0,0,67,0,0)$ & $2$ & $5$ & $23$ & $2$ & $(8,40,62,1,3,26)$\\
\bottomrule
\end{longtable}

\endgroup

\paragraph{Certificate supplement: all seventy-eight common-factor witnesses.}
For each saved pair $(\eta,r)$ below, the common factor is $x^\eta$ in
\[
 x^\eta\bigl(\lambda x_1^{10}x_4^{20}x_5^{10}
               +\mu x_0^{10}x_2^{30}\bigr),
\]
and $r$ is its quadratic-period witness. Coordinate indices start at zero;
the column $j$ is the wrap coordinate, corresponding to the added
$70\mathbf e_j$ in the reduction. These are finite certificates for the
specified family, not all possible tangents.\src{C3}

\begin{longtable}{rccr}
\toprule
No. & $\eta$ & $r$ & $j$\\
\midrule
\endhead
$1$ & $(3,0,0,15,5,7)$ & $(35,30,4,12,1,58)$ & $5$\\
$2$ & $(3,0,0,16,4,7)$ & $(35,30,4,10,3,58)$ & $5$\\
$3$ & $(3,0,0,16,5,6)$ & $(35,30,4,10,1,60)$ & $5$\\
$4$ & $(3,0,0,17,3,7)$ & $(35,30,4,8,5,58)$ & $5$\\
$5$ & $(3,0,0,17,4,6)$ & $(35,30,4,8,3,60)$ & $5$\\
$6$ & $(3,0,0,18,3,6)$ & $(35,30,4,6,5,60)$ & $5$\\
$7$ & $(3,0,1,14,5,7)$ & $(35,30,2,14,1,58)$ & $5$\\
$8$ & $(3,0,1,15,4,7)$ & $(35,30,2,12,3,58)$ & $5$\\
$9$ & $(3,0,1,15,5,6)$ & $(35,30,2,12,1,60)$ & $5$\\
$10$ & $(3,0,1,16,3,7)$ & $(35,30,2,10,5,58)$ & $5$\\
$11$ & $(3,0,1,16,4,6)$ & $(35,30,2,10,3,60)$ & $5$\\
$12$ & $(3,0,1,17,3,6)$ & $(35,30,2,8,5,60)$ & $5$\\
$13$ & $(3,1,0,14,5,7)$ & $(35,28,4,14,1,58)$ & $5$\\
$14$ & $(3,1,0,15,4,7)$ & $(35,28,4,12,3,58)$ & $5$\\
$15$ & $(3,1,0,15,5,6)$ & $(35,28,4,12,1,60)$ & $5$\\
$16$ & $(3,1,0,16,3,7)$ & $(35,28,4,10,5,58)$ & $5$\\
$17$ & $(3,1,0,16,4,6)$ & $(35,28,4,10,3,60)$ & $5$\\
$18$ & $(3,1,0,17,3,6)$ & $(35,28,4,8,5,60)$ & $5$\\
$19$ & $(3,1,1,14,4,7)$ & $(35,28,2,14,3,58)$ & $5$\\
$20$ & $(3,1,1,14,5,6)$ & $(35,28,2,14,1,60)$ & $5$\\
$21$ & $(3,1,1,15,3,7)$ & $(35,28,2,12,5,58)$ & $5$\\
$22$ & $(3,1,1,15,4,6)$ & $(35,28,2,12,3,60)$ & $5$\\
$23$ & $(3,1,1,16,3,6)$ & $(35,28,2,10,5,60)$ & $5$\\
$24$ & $(3,2,0,14,4,7)$ & $(35,26,4,14,3,58)$ & $5$\\
$25$ & $(3,2,0,14,5,6)$ & $(35,26,4,14,1,60)$ & $5$\\
$26$ & $(3,2,0,15,3,7)$ & $(35,26,4,12,5,58)$ & $5$\\
$27$ & $(3,2,0,15,4,6)$ & $(35,26,4,12,3,60)$ & $5$\\
$28$ & $(3,2,0,16,3,6)$ & $(35,26,4,10,5,60)$ & $5$\\
$29$ & $(3,2,1,14,3,7)$ & $(35,26,2,14,5,58)$ & $5$\\
$30$ & $(3,2,1,14,4,6)$ & $(35,26,2,14,3,60)$ & $5$\\
$31$ & $(3,2,1,15,3,6)$ & $(35,26,2,12,5,60)$ & $5$\\
$32$ & $(3,3,0,14,3,7)$ & $(35,24,4,14,5,58)$ & $5$\\
$33$ & $(3,3,0,14,4,6)$ & $(35,24,4,14,3,60)$ & $5$\\
$34$ & $(3,3,0,15,3,6)$ & $(35,24,4,12,5,60)$ & $5$\\
$35$ & $(3,3,1,14,3,6)$ & $(35,24,2,14,5,60)$ & $5$\\
$36$ & $(3,4,0,14,3,6)$ & $(35,22,4,14,5,60)$ & $5$\\
$37$ & $(4,0,0,14,5,7)$ & $(33,30,4,14,1,58)$ & $5$\\
$38$ & $(4,0,0,15,4,7)$ & $(33,30,4,12,3,58)$ & $5$\\
$39$ & $(4,0,0,15,5,6)$ & $(33,30,4,12,1,60)$ & $5$\\
$40$ & $(4,0,0,16,3,7)$ & $(33,30,4,10,5,58)$ & $5$\\
$41$ & $(4,0,0,16,4,6)$ & $(33,30,4,10,3,60)$ & $5$\\
$42$ & $(4,0,0,17,3,6)$ & $(33,30,4,8,5,60)$ & $5$\\
$43$ & $(4,0,1,14,4,7)$ & $(33,30,2,14,3,58)$ & $5$\\
$44$ & $(4,0,1,14,5,6)$ & $(33,30,2,14,1,60)$ & $5$\\
$45$ & $(4,0,1,15,3,7)$ & $(33,30,2,12,5,58)$ & $5$\\
$46$ & $(4,0,1,15,4,6)$ & $(33,30,2,12,3,60)$ & $5$\\
$47$ & $(4,0,1,16,3,6)$ & $(33,30,2,10,5,60)$ & $5$\\
$48$ & $(4,1,0,14,4,7)$ & $(33,28,4,14,3,58)$ & $5$\\
$49$ & $(4,1,0,14,5,6)$ & $(33,28,4,14,1,60)$ & $5$\\
$50$ & $(4,1,0,15,3,7)$ & $(33,28,4,12,5,58)$ & $5$\\
$51$ & $(4,1,0,15,4,6)$ & $(33,28,4,12,3,60)$ & $5$\\
$52$ & $(4,1,0,16,3,6)$ & $(33,28,4,10,5,60)$ & $5$\\
$53$ & $(4,1,1,14,3,7)$ & $(33,28,2,14,5,58)$ & $5$\\
$54$ & $(4,1,1,14,4,6)$ & $(33,28,2,14,3,60)$ & $5$\\
$55$ & $(4,1,1,15,3,6)$ & $(33,28,2,12,5,60)$ & $5$\\
$56$ & $(4,2,0,14,3,7)$ & $(33,26,4,14,5,58)$ & $5$\\
$57$ & $(4,2,0,14,4,6)$ & $(33,26,4,14,3,60)$ & $5$\\
$58$ & $(4,2,0,15,3,6)$ & $(33,26,4,12,5,60)$ & $5$\\
$59$ & $(4,2,1,14,3,6)$ & $(33,26,2,14,5,60)$ & $5$\\
$60$ & $(4,3,0,14,3,6)$ & $(33,24,4,14,5,60)$ & $5$\\
$61$ & $(5,0,0,14,4,7)$ & $(31,30,4,14,3,58)$ & $5$\\
$62$ & $(5,0,0,14,5,6)$ & $(31,30,4,14,1,60)$ & $5$\\
$63$ & $(5,0,0,15,3,7)$ & $(31,30,4,12,5,58)$ & $5$\\
$64$ & $(5,0,0,15,4,6)$ & $(31,30,4,12,3,60)$ & $5$\\
$65$ & $(5,0,0,16,3,6)$ & $(31,30,4,10,5,60)$ & $5$\\
$66$ & $(5,0,1,14,3,7)$ & $(31,30,2,14,5,58)$ & $5$\\
$67$ & $(5,0,1,14,4,6)$ & $(31,30,2,14,3,60)$ & $5$\\
$68$ & $(5,0,1,15,3,6)$ & $(31,30,2,12,5,60)$ & $5$\\
$69$ & $(5,1,0,14,3,7)$ & $(31,28,4,14,5,58)$ & $5$\\
$70$ & $(5,1,0,14,4,6)$ & $(31,28,4,14,3,60)$ & $5$\\
$71$ & $(5,1,0,15,3,6)$ & $(31,28,4,12,5,60)$ & $5$\\
$72$ & $(5,1,1,14,3,6)$ & $(31,28,2,14,5,60)$ & $5$\\
$73$ & $(5,2,0,14,3,6)$ & $(31,26,4,14,5,60)$ & $5$\\
$74$ & $(6,0,0,14,3,7)$ & $(29,30,4,14,5,58)$ & $5$\\
$75$ & $(6,0,0,14,4,6)$ & $(29,30,4,14,3,60)$ & $5$\\
$76$ & $(6,0,0,15,3,6)$ & $(29,30,4,12,5,60)$ & $5$\\
$77$ & $(6,0,1,14,3,6)$ & $(29,30,2,14,5,60)$ & $5$\\
$78$ & $(6,1,0,14,3,6)$ & $(29,28,4,14,5,60)$ & $5$\\
\bottomrule
\end{longtable}

\paragraph{Certificate supplement: all one hundred formal cover checks.}
At every listed multiplier $f$, $m=70f$ and $\nu(fa)$ is outside the
specified generator span modulo $2$. Only the mod-$2$ rank is tabulated;
the obstruction flag is true in every row. The four blocks of columns below
are a compact presentation of the complete $100$-row JSON record. No claim
is made beyond these finite levels.\src{C4}

\begingroup\small\setlength{\tabcolsep}{3pt}

\begin{longtable}{rrr@{\hspace{8pt}}rrr@{\hspace{8pt}}rrr@{\hspace{8pt}}rrr}
\toprule
$f$ & $m$ & Rank & $f$ & $m$ & Rank & $f$ & $m$ & Rank & $f$ & $m$ & Rank\\
\midrule
\endhead
$1$ & $70$ & $55$ & $26$ & $1820$ & $1523$ & $51$ & $3570$ & $3177$ & $76$ & $5320$ & $4447$\\
$2$ & $140$ & $111$ & $27$ & $1890$ & $1669$ & $52$ & $3640$ & $3055$ & $77$ & $5390$ & $4545$\\
$3$ & $210$ & $181$ & $28$ & $1960$ & $1619$ & $53$ & $3710$ & $3081$ & $78$ & $5460$ & $4867$\\
$4$ & $280$ & $227$ & $29$ & $2030$ & $1689$ & $54$ & $3780$ & $3339$ & $79$ & $5530$ & $4589$\\
$5$ & $350$ & $287$ & $30$ & $2100$ & $1851$ & $55$ & $3850$ & $3245$ & $80$ & $5600$ & $4635$\\
$6$ & $420$ & $363$ & $31$ & $2170$ & $1805$ & $56$ & $3920$ & $3243$ & $81$ & $5670$ & $5017$\\
$7$ & $490$ & $403$ & $32$ & $2240$ & $1851$ & $57$ & $3990$ & $3549$ & $82$ & $5740$ & $4771$\\
$8$ & $560$ & $459$ & $33$ & $2310$ & $2061$ & $58$ & $4060$ & $3379$ & $83$ & $5810$ & $4821$\\
$9$ & $630$ & $553$ & $34$ & $2380$ & $1987$ & $59$ & $4130$ & $3429$ & $84$ & $5880$ & $5199$\\
$10$ & $700$ & $575$ & $35$ & $2450$ & $2027$ & $60$ & $4200$ & $3711$ & $85$ & $5950$ & $4985$\\
$11$ & $770$ & $645$ & $36$ & $2520$ & $2223$ & $61$ & $4270$ & $3545$ & $86$ & $6020$ & $5003$\\
$12$ & $840$ & $735$ & $37$ & $2590$ & $2153$ & $62$ & $4340$ & $3611$ & $87$ & $6090$ & $5409$\\
$13$ & $910$ & $761$ & $38$ & $2660$ & $2219$ & $63$ & $4410$ & $3901$ & $88$ & $6160$ & $5191$\\
$14$ & $980$ & $807$ & $39$ & $2730$ & $2433$ & $64$ & $4480$ & $3707$ & $89$ & $6230$ & $5169$\\
$15$ & $1050$ & $925$ & $40$ & $2800$ & $2315$ & $65$ & $4550$ & $3825$ & $90$ & $6300$ & $5571$\\
$16$ & $1120$ & $923$ & $41$ & $2870$ & $2385$ & $66$ & $4620$ & $4123$ & $91$ & $6370$ & $5357$\\
$17$ & $1190$ & $993$ & $42$ & $2940$ & $2595$ & $67$ & $4690$ & $3893$ & $92$ & $6440$ & $5375$\\
$18$ & $1260$ & $1107$ & $43$ & $3010$ & $2501$ & $68$ & $4760$ & $3983$ & $93$ & $6510$ & $5781$\\
$19$ & $1330$ & $1109$ & $44$ & $3080$ & $2591$ & $69$ & $4830$ & $4293$ & $94$ & $6580$ & $5467$\\
$20$ & $1400$ & $1155$ & $45$ & $3150$ & $2785$ & $70$ & $4900$ & $4055$ & $95$ & $6650$ & $5565$\\
$21$ & $1470$ & $1297$ & $46$ & $3220$ & $2683$ & $71$ & $4970$ & $4125$ & $96$ & $6720$ & $5943$\\
$22$ & $1540$ & $1291$ & $47$ & $3290$ & $2733$ & $72$ & $5040$ & $4455$ & $97$ & $6790$ & $5633$\\
$23$ & $1610$ & $1341$ & $48$ & $3360$ & $2967$ & $73$ & $5110$ & $4241$ & $98$ & $6860$ & $5679$\\
$24$ & $1680$ & $1479$ & $49$ & $3430$ & $2839$ & $74$ & $5180$ & $4307$ & $99$ & $6930$ & $6201$\\
$25$ & $1750$ & $1447$ & $50$ & $3500$ & $2895$ & $75$ & $5250$ & $4645$ & $100$ & $7000$ & $5795$\\
\bottomrule
\end{longtable}

\endgroup

\paragraph{Complete source archives, code, and output provenance.}
Four original ZIP files were supplied in the conversation. Each is retained
byte-for-byte, with its original scripts, notes, finite outputs, and
certificates. The two separately supplied README files for the continuation
and third pass match their corresponding archive entries byte-for-byte.
The inventories below include every archive member other than directory
entries. Filenames identify computational artifacts, not additional
mathematical assumptions.\src{C1}\src{C2}\src{C3}\src{C4}

\textbf{hodge\_counterexample\_search.zip}\src{C1}

\begin{longtable}{@{}p{0.77\textwidth}r@{}}
\toprule
Archive member & Bytes\\
\midrule
\endhead
\path{README.md} & $6,922$\\
\path{search.cpp} & $2,208$\\
\path{independent_verify.py} & $3,361$\\
\path{lattice_check.py} & $956$\\
\path{jacobi_check.py} & $2,379$\\
\path{residual70.txt} & $107$\\
\path{verification_results.json} & $4,243$\\
\path{jacobi_results.json} & $697$\\
\bottomrule
\end{longtable}

\textit{SHA-256 of the original ZIP:}\par\nolinkurl{fb11e1493a13a149d0331ccb5527361beeaa6770c2201cd38983d5c9d859de26}

\textbf{hodge\_push\_further.zip}\src{C2}

\begin{longtable}{@{}p{0.77\textwidth}r@{}}
\toprule
Archive member & Bytes\\
\midrule
\endhead
\path{all_monomials.cpp} & $2,033$\\
\path{README.md} & $14,729$\\
\path{jacobi_results.json} & $882$\\
\path{tangent_certificate.py} & $1,729$\\
\path{high_order_directions.csv} & $2,518$\\
\path{single_class_support.py} & $1,010$\\
\path{monomial_census.txt} & $559$\\
\path{tangent_certificate.json} & $388$\\
\path{verify_deformations.py} & $2,921$\\
\path{deformation_results.json} & $813$\\
\path{single_class_support.json} & $474$\\
\path{jacobi_exact.py} & $2,445$\\
\bottomrule
\end{longtable}

\textit{SHA-256 of the original ZIP:}\par\nolinkurl{53937c5a7c3cadef6ebf3a97b4c589a96688a36ac9ae5bf52c4a0e3f27dfba04}

\textbf{hodge\_third\_pass.zip}\src{C3}

\begin{longtable}{@{}p{0.77\textwidth}r@{}}
\toprule
Archive member & Bytes\\
\midrule
\endhead
\path{README.md} & $10,601$\\
\path{verify.py} & $8,485$\\
\path{verification_results.json} & $24,344$\\
\path{pair_family.cpp} & $1,623$\\
\bottomrule
\end{longtable}

\textit{SHA-256 of the original ZIP:}\par\nolinkurl{acf24e0abed9ab3a2a949ace9d6d1ad3d8ce89ab63d84597abf7943002c6ebbb}

\textbf{hodge\_final\_attempt.zip}\src{C4}

\begin{longtable}{@{}p{0.77\textwidth}r@{}}
\toprule
Archive member & Bytes\\
\midrule
\endhead
\path{README.md} & $14,632$\\
\path{verify.py} & $8,623$\\
\path{all_gamma_certificates.json} & $16,373$\\
\path{lift_checks.json} & $11,538$\\
\path{verification_results.json} & $18,316$\\
\path{SHA256.json} & $459$\\
\bottomrule
\end{longtable}

\textit{SHA-256 of the original ZIP:}\par\nolinkurl{b3ff229cd5ea3c533db8285d3e8086ed9db935cdd2ba97f658e4c5f08dd8e87b}

\paragraph{How the standalone \texttt{.tex} preserves the original progress.}
The four ZIP archives are also embedded, byte-for-byte in Base64-encoded
comment blocks, after this file's \verb|\end{document}|. They do not affect
compilation and require no special LaTeX package. The readable mathematical
record above is self-contained; the embedded payload preserves the original
code, full Gamma multiplication certificates, all JSON rows, the detailed
README arguments, and discovery-era outputs without silently rewriting them.
The companion bundle supplies \texttt{extract\_hodge\_archives.py} to recover
those four ZIP files from the single \texttt{.tex}:
\begin{Code}
python extract_hodge_archives.py hodge_research_complete.tex recovered_archives
\end{Code}
The extractor verifies the recorded SHA-256 digest before writing each file.
It never executes the recovered research programs. The following compact
standard-library extractor is included here so that archive recovery does
not depend on obtaining the companion bundle:
\begin{Code}
from pathlib import Path
import base64, hashlib, re, sys

source, dest = Path(sys.argv[1]), Path(sys.argv[2])
dest.mkdir(parents=True, exist_ok=True)
pattern = re.compile(
    r"^% BEGIN_ARCHIVE ([\w.-]+) ([0-9a-f]{64})\n"
    r"((?:% [A-Za-z0-9+/=]+\n)+)% END_ARCHIVE\n", re.M)
for match in pattern.finditer(source.read_text(encoding="utf-8")):
    name, digest, body = match.groups()
    if Path(name).name != name:
        raise ValueError("Unsafe archive name")
    encoded = "".join(line[2:] for line in body.splitlines())
    data = base64.b64decode(encoded, validate=True)
    if hashlib.sha256(data).hexdigest() != digest:
        raise ValueError("Archive checksum mismatch")
    target = dest / name
    if target.exists() and target.read_bytes() != data:
        raise FileExistsError(target)
    target.write_bytes(data)
    print(target)
\end{Code}

\paragraph{Reproduction commands retained from the research rounds.}
Run the following only after extracting the relevant ZIP and entering its
research directory. Python requirements of the first two rounds are Python
3, NumPy, and SymPy; optional independent censuses require a C++17 compiler.
The third and fourth Python verifiers use the standard library. These
commands reproduce the specified finite checks, not a formal proof of the
geometric interpretation or the Hodge conjecture.
\begin{Code}
# hodge_counterexample_search/
python independent_verify.py
python jacobi_check.py
g++ -O3 -std=c++17 search.cpp -o search
./search 70 > residual70.txt
python lattice_check.py

# hodge_push_further/
python tangent_certificate.py
python verify_deformations.py
python jacobi_exact.py
python single_class_support.py
g++ -O3 -std=c++17 all_monomials.cpp -o all_monomials
./all_monomials > monomial_census.txt

# hodge_third_pass/
python verify.py
python verify.py --family
g++ -O2 -std=c++17 pair_family.cpp -o pair_family
./pair_family

# hodge_final_attempt/
python verify.py
python verify.py --covers
\end{Code}
Some programs rewrite their JSON reports in their own directory. The
original ZIPs should therefore be kept unchanged, with reruns performed in
separate working copies. Discovery used integer linear algebra and, for the
short Gamma certificate, mixed-integer optimization; the shipped final
verifier needs neither an optimizer nor floating-point arithmetic.

\paragraph{Export-time checks and what they do not certify.}
During preparation of this consolidated document, the original archives were
inspected, their hashes recorded, and the default core verifiers from
\texttt{hodge\_third\_pass} and \texttt{hodge\_final\_attempt} were rerun in
separate copies. Both reported that all their exact checks passed. The
third-pass rerun checked the original two-monomial support and quadratic
certificate, not its optional full family census. The final-pass rerun
checked the $24$ Gamma ratios, the trace normalization, the $32$
Jacobi--Kummer comparisons, and the divisor-enlarged degree-$70$ formal
obstruction, not the optional $100$-level rerun. Those extended census
results are retained from their supplied original reports.

The other original exhaustive programs were not rerun merely for this
export. No fresh literature survey or current-news verification was
performed. The supplied finite programs are not independent peer review of
the geometric arguments, and their passing is not a Lean proof. The
LaTeX source was separately compiled and checked as a document; this is a
formatting check, not a mathematical validation.

\paragraph{Verification and outcome.}
The concrete candidate is the smooth projective fourfold
\eqref{hodge:fermat}, with the rational block \eqref{hodge:block} and explicit
class \eqref{hodge:explicit-alpha}. The record retains the exact character
census, selected formal-lattice exclusions, finite-order Frobenius values,
monomial escape census, eighth-order period coefficient, simultaneous
finite-cover argument, uncorrectable second-order two-monomial obstruction,
Gamma certificates, Kummer equation and degree argument, $32$ prime
comparisons, and $100$ bounded formal-cover tests. The later reversed search
adds the explicit complete-intersection witness conditions
\eqref{hodge:witness-decomposition} and \eqref{hodge:witness-coefficient}.

\textbf{Outcome: explicit candidate, exact finite certificates, and identified
obstructions to several search routes; UNRESOLVED.} No non-algebraic rational
Hodge class has been established. The remaining counterexample claim is
\eqref{hodge:remaining-alpha}, equivalently
\eqref{hodge:remaining-block} for this block. The opposite constructive goal
is a genuine algebraic surface with nonzero projected class; no such
surface was produced. Nothing in the reported-breakthrough discussion is
used to close either gap.

\subsection{Sources}
The source list distinguishes the supplied authoritative statement,
established mathematical inputs cited in the conversation, original
research artifacts, and unverified reporting provenance. Bibliographic
details and public URLs from the earlier conversation and README files are
preserved; they were not independently re-retrieved for this export. In
particular, no new ``consulted'' date is invented for a public page.

\begin{itemize}
\item[S1]\hypertarget{src:S1}{}Clay Mathematics Institute,
\emph{The Millennium Prize Problems}.
\url{https://www.claymath.org/library/monographs/MPPc.pdf}.
\textit{Background formulation collection; linked in the supplied template.}

\item[E1]\hypertarget{src:E1}{}Pierre Deligne, \emph{The Hodge Conjecture},
official Clay problem description, supplied as \texttt{hodge.pdf}, five
pages. Section 1, especially page 2, is the exact rational smooth-projective
statement; Section 2 explains the integral and non-projective distinctions.
\url{https://www.claymath.org/wp-content/uploads/2022/06/hodge.pdf}.
\textit{The supplied PDF is authoritative for this export; the URL is a
public locator, not a newly checked status source.}

\item[S2]\hypertarget{src:S2}{}Clay Mathematics Institute,
\emph{Hodge Conjecture}, institutional overview.
\url{https://www.claymath.org/millennium/hodge-conjecture/}.
\textit{Referenced in the earlier status discussion; no present-day status
claim is inferred from that earlier reference.}

\item[E2]\hypertarget{src:E2}{}Genival da Silva Jr.,
\emph{Notes on the Hodge Conjecture for Fermat Varieties},
Experimental Results 2 (2021), e22; arXiv:2101.04739.
\url{https://arxiv.org/html/2101.04739}.
\textit{Fermat character decomposition, rational Hodge criterion, and
cycle-construction framework cited by the research archives.}

\item[E3]\hypertarget{src:E3}{}Noboru Aoki,
\emph{Some new algebraic cycles on Fermat varieties},
Journal of the Mathematical Society of Japan 39 (1987), 385--396; and
\emph{Some remarks on the Hodge conjecture for abelian varieties of Fermat
type}, Commentarii Mathematici Universitatis Sancti Pauli 49 (2000),
177--194.
\textit{Bibliographic entries as retained in the archives. The implemented
standard-generator list was accessed there through the later source [E4];
no completeness theorem for all algebraic cycles is assumed.}

\item[E4]\hypertarget{src:E4}{}Rifat Jumagulov,
\emph{The Hodge conjecture for Fermat fourfolds of odd degree at most 199},
arXiv:2608.18134v1 (2026), especially Appendices B and D, as cited in the
prior conversation and archives.
\url{https://arxiv.org/html/2608.18134v1}.
\textit{Inherited attribution for a search direction and formal-generator
conventions; the preprint's existence, contents, and reported census were
not independently reverified while exporting. It is not cited as a proof
that the selected block is non-algebraic.}

\item[E5]\hypertarget{src:E5}{}Andr\'e Weil,
\emph{Numbers of solutions of equations in finite fields},
Bulletin of the American Mathematical Society 55 (1949), 497--508.
\url{https://download.uni-mainz.de/mathematik/Algebraische%20Geometrie/Lehre/WS23.Padische.Weil.nf.pdf}.
\textit{Diagonal equations and Jacobi sums; public copy linked in the first
research response.}

\item[E6]\hypertarget{src:E6}{}Mark Watkins,
\emph{Jacobi sums and Grossencharacters}, Publications math\'ematiques de
Besan\c{c}on (2018), 111--122.
\url{https://www.numdam.org/item/10.5802/pmb.25.pdf}.
\textit{Hecke characters, infinity types, and arithmetic conventions, as
cited in the continuation archive.}

\item[E7]\hypertarget{src:E7}{}Pierre Lairez,
\emph{Computing periods of rational integrals}, arXiv:1404.5069,
Section 1.1.
\url{https://arxiv.org/html/1404.5069v3}.
\textit{Griffiths--Dwork reduction identity used in the deformation
calculations.}

\item[E8]\hypertarget{src:E8}{}Hossein Movasati,
\emph{Why should one compute periods of algebraic cycles?},
arXiv:1602.06607, Theorem 5 as cited in the archive.
\url{https://arxiv.org/html/1602.06607}.
\textit{Period-matrix description of Hodge tangent spaces.}

\item[E9]\hypertarget{src:E9}{}Jorge Duque Franco and Roberto Villaflor Loyola,
\emph{On fake linear cycles inside Fermat varieties},
Algebra \& Number Theory 17 (2023), no. 10, 1847--1865;
arXiv:2112.14818v2, Proposition 3.1 and the preceding discussion of spanning
vanishing cycles, Proposition 3.2, and Sections 4 and 6 as cited in the
archives.
\url{https://arxiv.org/html/2112.14818v2}.
\textit{Explicit Fermat periods, Hodge tangent ideals, and second-order
framework. The specific two-monomial obstruction is attributed to the
recorded investigation, not to this paper.}

\item[E10]\hypertarget{src:E10}{}NIST Digital Library of Mathematical
Functions, Section 5.5, especially equation 5.5.3 (Euler reflection) and
5.5.6 (Gauss multiplication).
\url{https://dlmf.nist.gov/5.5}.
\textit{Established special-function identities; the combinations and
finite exponent certificates are supplied by the research archives.}

\item[E11]\hypertarget{src:E11}{}Pierre Deligne, notes by J. S. Milne,
\emph{Hodge Cycles on Abelian Varieties}, Lecture Notes in Mathematics 900
(1982), Section 7, especially the Fermat discussion and Theorem 7.15;
updated typeset version cited in the archive, pages 45 and 53--55.
\url{https://www.jmilne.org/math/Documents/Deligne82.pdf}.
\textit{General normalized Fermat-period algebraicity and period--Galois
comparison, not a proof of algebraicity or non-algebraicity of the
particular cycle candidate.}

\item[E12]\hypertarget{src:E12}{}Roberto Villaflor Loyola,
complete-intersection cycle-class and Jacobian-determinant formula cited in
the last mathematical response, arXiv:1812.03964.
\url{https://arxiv.org/html/1812.03964}.
\textit{Reference preserved at the level supplied by the conversation;
no additional publication details or theorem number are invented here.}

\item[C1]\hypertarget{src:C1}{}Original first-pass research archive,
\texttt{hodge\_counterexample\_search.zip}.
\textit{Contains the initial README, exact sextuple census, independent
Python verification, formal-lattice checks, Jacobi-sum checks, six residual
representatives, and saved JSON outputs. Supplied in the conversation and
embedded unchanged after the end of this LaTeX document.}

\item[C2]\hypertarget{src:C2}{}Original continuation archive,
\texttt{hodge\_push\_further.zip}.
\textit{Contains the complete monomial census, degree-$71$ tangent
certificate, eighth-order calculation, exact order-$70$ Jacobi data,
single-class support pattern, geometric notes, and outputs. The separately
supplied README matches its archived copy.}

\item[C3]\hypertarget{src:C3}{}Original third-pass archive,
\texttt{hodge\_third\_pass.zip}.
\textit{Contains the second-order obstruction argument, exact Gamma
certificate, six quadratic rows, independent family-census code, and the
$78$ common-factor witness rows. The separately supplied README matches its
archived copy.}

\item[C4]\hypertarget{src:C4}{}Original final-pass archive,
\texttt{hodge\_final\_attempt.zip}.
\textit{Contains the explicit rational class, period normalization,
all inverse-pair Gamma certificates, the Kummer equation and valuation
argument, $32$ prime comparisons, $100$ lift checks, source hashes, and
standard-library verifier.}

\item[C5]\hypertarget{src:C5}{}The supplied conversation preceding this
export.
\textit{Source for the sequence of user requests, the inherited progress
reports, the rumor discussion and its as-of markers, the existence
heuristic, and the final complete-intersection witness strategy. The last
strategy has no separately supplied computational archive.}

\item[T1]\hypertarget{src:T1}{}User-supplied formatting example,
\texttt{Pasted text(20260927-133058).txt}, containing a standalone BSD
research-attempt LaTeX document.
\textit{Used only for document structure and typography: one problem
section, definitions, short English statement, research attempt, and
sources. Its BSD-specific statement, references, rank, and claims are not
reused as Hodge facts.}

\item[N1]\hypertarget{src:N1}{}MTSlive post supplied by the user.
\url{https://x.com/MTSlive/status/2100626270774263813}.
\textit{User-provided reporting claim about OpenAI and Hodge, quoted in the
provenance paragraph; not a mathematical argument or an independently
verified announcement in this export.}

\item[N2]\hypertarget{src:N2}{}The Information article URL linked by the
earlier assistant:
\url{https://www.theinformation.com/newsletters/ai-agenda/openai-close-solving-another-millennium-prize-math-problem}.
\textit{The prior reply described the body as subscriber-only and
attributed a public summary to Stephanie Palazzolo. The supplied materials
do not include that article body or a separate permalink for the reporter's
statement. Neither is treated as proof of a mathematical result.}

\item[N3]\hypertarget{src:N3}{}OpenAI public research index, URL cited in
the earlier news response.
\url{https://openai.com/research/index/}.
\textit{Retained solely to identify the earlier assertion about what that
response said it found. No current index contents, Hodge solution, or
separate Navier--Stokes announcement are verified or endorsed by this
export.}
\end{itemize}
\end{document}

%
